% why-notes: true
% expositor: true

\section{环的定义与例子}\label{sec:ring-definition-examples}

\subsection{环的定义}\label{sec:ring-definition}

上一章讲域时，我们把"可以做四则运算"这件事公理化了。
但很多自然出现的代数对象并没有那么好：矩阵可以加、可以乘，却不满足乘法交换；整数可以加、可以乘，却不是每个非零整数都有乘法逆元。
如果仍然坚持域的定义，这些对象都会被排除在外。

环就是在这个地方出现的。
粗略地说，环仍然有加法和乘法，但乘法的要求比域弱得多：它要满足结合律，并且要和加法通过分配律配合；至于乘法是否交换、是否有单位元、非零元素是否有逆元，都先不强求。

\paragraph*{公理化定义}

\begin{definition}[环]\label{def:ring}
设 $R$ 是一个非空集合，配备两种二元运算 $+$ 和 $\cdot$。
若满足以下条件，则称 $R$ 为一个\textbf{环}：

\begin{enumerate}
    \item \textbf{加法阿贝尔群}：$(R,+)$ 构成阿贝尔群。也就是说，对任意 $a,b,c\in R$，
    \begin{itemize}
        \item $(a+b)+c=a+(b+c)$；
        \item $a+b=b+a$；
        \item 存在加法单位元 $0\in R$，使得 $a+0=a$；
        \item 对任意 $a\in R$，存在加法逆元 $-a\in R$，使得 $a+(-a)=0$。
    \end{itemize}
    \item \textbf{乘法结合律}：对任意 $a,b,c\in R$，
    \[
    (ab)c=a(bc).
    \]
    \item \textbf{左右分配律}：对任意 $a,b,c\in R$，
    \[
    a(b+c)=ab+ac,\qquad (a+b)c=ac+bc.
    \]
\end{enumerate}
\end{definition}

这里的乘法通常省略点号，直接写成 $ab$。
注意，定义 \ref{def:ring} 并没有要求乘法单位元，也没有要求乘法交换。
如果环 $R$ 中存在元素 $1_R$，使得
\[
1_Ra=a1_R=a,\qquad a\in R,
\]
则称 $R$ 是\textbf{有单位元的环}，也常说\textbf{含幺环}。
如果对任意 $a,b\in R$ 都有 $ab=ba$，则称 $R$ 是\textbf{交换环}。
所以"有单位元"和"交换"是两件不同的事；同时满足二者时，才说它是交换含幺环。

环的定义和域的定义相比，真正放松的是乘法部分。
域要求非零元素在乘法下构成阿贝尔群；环只要求乘法结合，并通过分配律与加法相容。
这正好解释了为什么环的例子会比域丰富得多。

\subsection{环的基本例子}\label{sec:ring-basic-examples}

判断一个对象是不是环，不能只看它有加法和乘法，还要看这些运算是否满足定义中的相容条件。
先从线性代数里最熟悉的对象开始。

\paragraph*{例 1：线性变换环和矩阵环}
设 $V$ 是数域 $F$ 上的有限维线性空间。
从 $V$ 到自身的所有线性变换记为
\[
\operatorname{End}_F(V).
\]
两个线性变换可以相加，也可以复合。
加法按点定义：
\[
(S+T)(v)=S(v)+T(v),\qquad v\in V.
\]
因为线性变换之和仍是线性变换，零映射是加法单位元，$-T$ 仍是线性变换，所以 $\operatorname{End}_F(V)$ 在加法下构成阿贝尔群。

乘法取为复合：
\[
ST=S\circ T.
\]
复合仍是线性变换，且映射复合满足结合律：
\[
(RS)T=R(ST).
\]
分配律来自线性变换加法和复合的相容性。对任意 $v\in V$，
\[
((S+T)R)(v)=(S+T)(R(v))=S(R(v))+T(R(v)),
\]
所以
\[
(S+T)R=SR+TR.
\]
同理，
\[
(R(S+T))(v)=R(S(v)+T(v))=R(S(v))+R(T(v)),
\]
因此 $R(S+T)=RS+RT$。
三条公理都已满足，所以 $\operatorname{End}_F(V)$ 是一个环。
它的乘法单位元就是恒等变换 $\operatorname{id}_V$。

如果给定 $V$ 的一组基，每个线性变换都可以表示成一个 $n$ 阶矩阵。
于是
\[
\operatorname{End}_F(V)\cong M_n(F).
\]
这说明 $n$ 阶矩阵全体在矩阵加法和矩阵乘法下构成环；矩阵加法、乘法的环公理可以看作线性变换环在选定基后的坐标表示。
当 $n\geq 2$ 时，这个环一般不交换。
这类例子最适合提醒我们：乘法不交换并不妨碍它成为环。

\paragraph*{例 2：$m\mathbb Z$}
设 $m$ 是一个正整数，记
\[
m\mathbb Z=\{mk:k\in\mathbb Z\}.
\]
它在通常的整数加法和乘法下构成一个环。
先看加法。若 $ma,mb\in m\mathbb Z$，则
\[
ma+mb=m(a+b)\in m\mathbb Z,
\]
零元 $0=m\cdot0$ 属于 $m\mathbb Z$，而 $ma$ 的加法逆元是
\[
-ma=m(-a)\in m\mathbb Z.
\]
加法结合律和交换律继承自整数加法，所以 $(m\mathbb Z,+)$ 是阿贝尔群。

乘法也封闭，因为
\[
(ma)(mb)=m(mab)\in m\mathbb Z.
\]
乘法结合律和左右分配律都继承自整数环 $\mathbb Z$。
因此 $m\mathbb Z$ 满足环定义。
当 $m>1$ 时，这个环没有乘法单位元。
比如若 $e\in m\mathbb Z$ 是单位元，则对 $m\in m\mathbb Z$ 应有 $em=m$，在整数中推出 $e=1$，但 $1\notin m\mathbb Z$。
这说明环的定义确实不能强行要求单位元，否则这类自然的例子就会被排除。

\paragraph*{例 3：高斯整数和二次整数}
高斯整数环定义为
\[
\mathbb Z[i]=\{a+bi:a,b\in\mathbb Z\}.
\]
它的加法和乘法沿用复数中的加法和乘法。
加法封闭是因为
\[
(a+bi)+(c+di)=(a+c)+(b+d)i.
\]
加法单位元是 $0=0+0i$，加法逆元是
\[
-(a+bi)=(-a)+(-b)i.
\]
乘法封闭则用到 $i^2=-1$：
\[
(a+bi)(c+di)=(ac-bd)+(ad+bc)i.
\]
右边系数仍在 $\mathbb Z$ 中。
结合律、交换律和分配律都继承自复数的运算。
因此 $\mathbb Z[i]$ 是一个交换含幺环。

类似地，若 $D$ 是无平方因子的整数，可以考虑
\[
\mathbb Z[\sqrt D]=\{a+b\sqrt D:a,b\in\mathbb Z\}.
\]
加法封闭和加法逆元的验证与 $\mathbb Z[i]$ 完全一样。
乘法封闭来自
\[
(a+b\sqrt D)(c+d\sqrt D)=(ac+bdD)+(ad+bc)\sqrt D.
\]
结合律、交换律和分配律从所在的数域或复数域中继承下来，单位元是 $1=1+0\sqrt D$。
这里的方括号有一个常用含义：$\mathbb Z[\sqrt D]$ 表示由 $\mathbb Z$ 和 $\sqrt D$ 生成的环。
与此相对，圆括号常用来表示生成的域，例如 $\mathbb Q(\sqrt D)$。
两种记号的区别在于：环只要求对加减乘封闭，而域还要允许非零元素取逆。

\paragraph*{例 4：函数环}
设 $X$ 是非空集合，$R$ 是一个已知的环。
记
\[
R^X=\{f:X\to R\}
\]
为所有从 $X$ 到 $R$ 的函数组成的集合。
在逐点加法和逐点乘法下，
\[
(f+g)(x)=f(x)+g(x),\qquad (fg)(x)=f(x)g(x),
\]
$R^X$ 构成一个环。
验证时只要把每条等式都拿到点 $x\in X$ 上看。
例如加法结合律是
\[
((f+g)+h)(x)=(f(x)+g(x))+h(x)=f(x)+(g(x)+h(x))=(f+(g+h))(x).
\]
零函数 $x\mapsto0_R$ 是加法单位元，函数 $-f$ 由 $(-f)(x)=-f(x)$ 定义。
乘法结合律同样逐点继承自 $R$：
\[
((fg)h)(x)=(f(x)g(x))h(x)=f(x)(g(x)h(x))=(f(gh))(x).
\]
左右分配律也逐点验证，例如
\[
(f(g+h))(x)=f(x)(g(x)+h(x))=f(x)g(x)+f(x)h(x).
\]
如果 $R$ 有单位元，那么常值函数 $x\mapsto 1_R$ 就是 $R^X$ 的单位元；如果 $R$ 交换，那么 $R^X$ 也交换。
这个例子说明，给定一个环以后，可以通过"逐点运算"制造出新的环。

\paragraph*{例 5：群代数}
设 $G$ 是有限群，$F$ 是域。
群代数 $F[G]$ 的构造思路是：先把群元素当成一组形式基，得到一个 $F$-向量空间；再把群里的乘法扩展成这个向量空间上的乘法。

作为集合，
\[
F[G]=\left\{\sum_{g\in G}a_g g:a_g\in F\right\}.
\]
这里的 $g$ 不是被当作 $F$ 中的元素，而是被当作形式基向量。
因为 $G$ 有限，$F[G]$ 就是以 $G$ 为基的有限维 $F$-向量空间。
加法按同一个基向量的系数相加：
\[
\sum_{g\in G}a_g g+\sum_{g\in G}b_g g
=\sum_{g\in G}(a_g+b_g)g.
\]
零元是 $\sum_{g\in G}0g$，而 $\sum a_g g$ 的加法逆元是 $\sum (-a_g)g$。
所以 $(F[G],+)$ 是阿贝尔群，环定义的第一部分已经满足。

加法已经由向量空间保证；新的结构集中在乘法怎么定义。
乘法先在基元素上定义：两个基元素 $g,h\in G$ 相乘，结果就是群里的乘积 $gh$。
一般元素的乘法则由双线性扩展给出：
\[
\left(\sum_{g\in G}a_g g\right)
\left(\sum_{h\in G}b_h h\right)
=
\sum_{g,h\in G}a_gb_h(gh).
\]
这里一项 $a_gb_h(gh)$ 的来源是：从左边取 $a_g g$，从右边取 $b_h h$，系数在 $F$ 里相乘，基元素在 $G$ 里相乘。
由于 $gh$ 仍是 $G$ 中的元素，乘积仍是这些形式基的有限线性组合，所以乘法封闭。

下面用一个最小的非平凡例子来展示这种乘法。
设
\[
G=C_2=\{e,s\},\qquad s^2=e.
\]
则 $F[C_2]$ 中的元素形如 $ae+bs$。
两个元素相乘时，
\[
(ae+bs)(ce+ds)
=ac\,e+ad\,s+bc\,s+bd\,s^2
=(ac+bd)e+(ad+bc)s.
\]
这个例子显示，群代数乘法不是逐坐标相乘，而是先按群乘法合并基向量，再收集同类项。

加法阿贝尔群已经由向量空间结构给出。
乘法结合律的来源藏在基元素里：对 $g,h,k\in G$，
\[
(gh)k=g(hk)
\]
正是群乘法的结合律。
一般线性组合的结合律由这个基元素等式按双线性展开得到；每一项都在比较同一个三元组 $(g,h,k)$ 在 $(gh)k$ 和 $g(hk)$ 中的系数。

分配律来自乘法定义本身的双线性。
例如在 $\alpha(\beta+\gamma)$ 中，固定一对基元素 $(g,h)$，对应项的系数为 $a_g(b_h+c_h)$，由 $F$ 中的分配律拆成 $a_gb_h+a_gc_h$，于是这一项变成
\[
a_gb_h(gh)+a_gc_h(gh),
\]
正好分别出现在 $\alpha\beta$ 和 $\alpha\gamma$ 中。
右分配律同理。

若 $e$ 是群 $G$ 的单位元，那么 $1_F e$ 是 $F[G]$ 的乘法单位元。
原因是对每个基元素 $g$ 都有
\[
(1_F e)g=g(1_F e)=g,
\]
再按线性性扩展到任意元素。
因此 $F[G]$ 是一个含幺环。
它通常不是交换环：如果 $G$ 非交换，那么不同的群元素作为基向量线性无关，$gh\neq hg$ 会直接给出两个不同的基向量；若 $G$ 是交换群，则 $F[G]$ 是交换环。

还有一个常用视角，可以看清群代数乘法与函数环逐点乘法的区别。
因为 $G$ 是有限集，$F[G]$ 中的形式和
\[
\sum_{g\in G}a_g g
\]
等价于一个函数
\[
a:G\to F,\qquad g\mapsto a(g).
\]
这里 $a(g)$ 就是基向量 $g$ 前面的坐标。
不过在这个视角下，群代数的乘法不是逐点乘法，而是卷积乘法：
\[
(a*b)(k)=\sum_{gh=k}a(g)b(h),\qquad k\in G.
\]
固定一个 $k\in G$ 后，求和指标不是所有单个群元素，而是所有有序对 $(g,h)$，它们满足 $gh=k$。
每一对 $(g,h)$ 都给出一个贡献 $a(g)b(h)$：它来自第一项中基向量 $g$ 的系数和第二项中基向量 $h$ 的系数。
这个公式是前面乘法公式的"收集同类项"版本：所有满足 $gh=k$ 的乘积项都会贡献到基向量 $k$ 的系数上。
有限性保证了这个求和是有限和。

\paragraph*{例 6：路径代数}
路径代数和群代数很像，只是把"群元素"换成了 quiver（有向图，允许重边和环路）中的路径。
当前例子只需要最简单的有向图情形。
顶点看成长度为 $0$ 的路径，箭头看成长度为 $1$ 的路径，首尾能接上的路径可以拼成更长的路径。

设 $Q$ 是一个 quiver，$F$ 是域。
路径代数 $FQ$ 是以 $Q$ 中所有路径为基的 $F$-向量空间。
因此它的加法阿贝尔群结构来自向量空间加法。
乘法也先在基路径上定义。
本节采用从右往左读的乘法约定：
\[
pq=\text{先走 }q,\text{ 再走 }p.
\]
如果 $q$ 的终点等于 $p$ 的起点，就把两条路径拼接起来；如果接不上，就规定乘积为 $0$。
一般线性组合再按线性性展开：
\[
\left(\sum_p a_p p\right)\left(\sum_q b_q q\right)=\sum_{p,q}a_pb_q(pq)
\]
给出。
这里 $p,q$ 遍历 $Q$ 中所有路径。
一项 $a_pb_q(pq)$ 的意思是：先从右边线性组合中取路径 $q$，再从左边线性组合中取路径 $p$；系数相乘为 $a_pb_q$，路径部分按约定先走 $q$ 再走 $p$。
如果 $q$ 的终点和 $p$ 的起点接不上，那么这一项中的 $pq$ 就是 $0$，不会贡献新的路径。
乘法封闭是因为拼接后的结果仍是路径，接不上时结果是 $0$。

结合律仍然先看基路径。
对三条基路径 $p,q,r$，$(pq)r$ 和 $p(qr)$ 非零当且仅当两个拼接处都匹配；这时二者都是按从右到左顺序拼出的同一条路径。
只要有一个拼接处不匹配，两边都会等于 $0$。
因此基路径上的乘法满足结合律。
一般线性组合由基路径的结果线性扩展而来。
分配律也一样：路径乘积 $pq$ 不变，只是系数里的 $a_p(b_q+c_q)$ 被拆成 $a_pb_q+a_pc_q$。
于是路径代数的环公理归结为两件事：路径拼接是结合的，系数域 $F$ 中的加法和乘法满足分配律。

先看一个小例子。
设 quiver（有向图）$Q$ 是
\[
1\xrightarrow{\alpha}2.
\]
路径只有
\[
e_1,\ e_2,\ \alpha,
\]
其中 $e_1,e_2$ 是长度为 $0$ 的路径。
因为 $\alpha:1\to2$，按从右往左读的约定，
\[
e_2\alpha=\alpha,\qquad \alpha e_1=\alpha.
\]
而
\[
e_1\alpha=0,\qquad \alpha e_2=0,\qquad \alpha^2=0,
\]
因为这些乘积对应的首尾都接不上。
所以这个例子中的路径代数就是
\[
FQ=\{ae_1+be_2+c\alpha:a,b,c\in F\},
\]
乘法完全由上面这些基路径规则和分配律决定。

再看一个稍复杂的路径代数。
\[
1\xrightarrow{\alpha}2\xrightarrow{\beta}3.
\]
若采用从右往左读乘法的约定，则先走 $\alpha$ 再走 $\beta$ 的路径写作 $\beta\alpha$，所以 $\beta\alpha\neq0$；而 $\alpha\beta$ 表示先走 $\beta$ 再走 $\alpha$，首尾接不上，因此为 $0$。
如果 $Q$ 的顶点集合有限，那么长度为 $0$ 的路径之和
\[
\sum_{i\in Q_0}e_i
\]
就是 $FQ$ 的单位元。
原因是任意路径 $p:i\to j$ 左边只会被终点处的 $e_j$ 接上，右边只会被起点处的 $e_i$ 接上，其他顶点项都因接不上而变成 $0$，于是
\[
\left(\sum_{v\in Q_0}e_v\right)p=p,\qquad
p\left(\sum_{v\in Q_0}e_v\right)=p.
\]
路径代数通常不交换。

\paragraph*{例 7：四元数环}
Hamilton 四元数可以写成
\[
\mathbb H=\{a+bi+cj+dk:a,b,c,d\in\mathbb R\},
\]
其中
\[
i^2=j^2=k^2=-1,\qquad ij=k,\quad jk=i,\quad ki=j,
\]
交换顺序则变号：
\[
ji=-k,\qquad kj=-i,\qquad ik=-j.
\]
加法就是 $\mathbb R^4$ 的逐坐标加法，所以加法阿贝尔群结构没有新困难。
乘法按基元乘法规则逐项展开。
封闭性来自这样一个事实：任意两个四元数相乘后仍能写成 $a+bi+cj+dk$ 的形式。
若记
\[
u_0=1,\qquad u_1=i,\qquad u_2=j,\qquad u_3=k,
\]
并写
\[
q=\sum_{\ell=0}^3x_\ell u_\ell,\qquad
q'=\sum_{m=0}^3y_m u_m,
\]
则乘法可以写成
\[
qq'=\sum_{\ell,m=0}^3x_\ell y_m(u_\ell u_m).
\]
这个双重求和一共有 $16$ 项。
每一项都来自第一个四元数中的一个基元 $x_\ell u_\ell$ 和第二个四元数中的一个基元 $y_m u_m$；系数相乘为 $x_\ell y_m$，基元乘积 $u_\ell u_m$ 再按乘法表化回 $\pm1,\pm i,\pm j,\pm k$。
例如 $\ell=1,m=2$ 时，得到的项是
\[
x_1y_2(u_1u_2)=x_1y_2(ij)=x_1y_2k.
\]
再取第三个四元数
$q''=\sum_{m=0}^3z_m u_m$。
逐项展开得到
\[
q(q'+q'')
=\sum_{\ell,m=0}^3x_\ell(y_m+z_m)(u_\ell u_m)
=\sum_{\ell,m=0}^3x_\ell y_m(u_\ell u_m)
 +\sum_{\ell,m=0}^3x_\ell z_m(u_\ell u_m)
=qq'+qq''.
\]
右边的分配律也由同样的逐项展开得到：
\[
(q+q')q''=qq''+q'q''.
\]
因此四元数乘法对实数是双线性的，左右分配律随之成立。

四元数乘法虽然不交换，但它仍然结合。
验证的一种方式是把四元数表示成实矩阵。例如令
\[
I=\begin{pmatrix}0&-1&0&0\\1&0&0&0\\0&0&0&-1\\0&0&1&0\end{pmatrix},\quad
J=\begin{pmatrix}0&0&-1&0\\0&0&0&1\\1&0&0&0\\0&-1&0&0\end{pmatrix},\quad
K=\begin{pmatrix}0&0&0&-1\\0&0&-1&0\\0&1&0&0\\1&0&0&0\end{pmatrix},
\]
则 $I^2=J^2=K^2=-\mathrm{id}$，$IJ=K$，$JK=I$，$KI=J$，与四元数乘法表一致。
四元数 $a+bi+cj+dk$ 对应矩阵 $aI_4+bI+cJ+dK$，矩阵乘法的结合律保证了四元数乘法也结合。
在这些乘法规则下，$\mathbb H$ 是一个有单位元的环，单位元是 $1$。
它的特点是每个非零四元数都有逆元。
若
\[
q=a+bi+cj+dk\neq0,
\]
定义它的共轭为
\[
\overline q=a-bi-cj-dk.
\]
直接按乘法规则展开可得
\[
q\overline q=\overline q q=a^2+b^2+c^2+d^2.
\]
因此
\[
q^{-1}=\frac{\overline q}{q\overline q}
=\frac{a-bi-cj-dk}{a^2+b^2+c^2+d^2}.
\]
分母是正实数，因为 $q\neq0$。
不过 $\mathbb H$ 的乘法不交换，例如 $ij=k$ 而 $ji=-k$。
所以它不是域（域要求乘法交换），而是一个除环（也称斜域）：每个非零元素都有逆元，但乘法不必交换。
$\mathbb H$ 是理解非交换乘法的重要例子。

\subsection{环的特征}\label{sec:ring-characteristic}

第一章讨论域的特征时，我们关心的是一个很朴素的问题：把乘法单位元 $1$ 反复相加，什么时候会回到 $0$。
这个问题并不只属于域。只要一个环有乘法单位元，同样可以问 $1_R$ 在加法群中生成的那一串元素
\[
0,\ 1_R,\ 2\cdot 1_R,\ 3\cdot 1_R,\ \ldots
\]
是否会回到 $0$。

本节只讨论含幺环的特征。原因是下面的定义要用到环中的乘法单位元 $1_R$。
前面例子中的矩阵环、函数环、群代数、有限顶点的路径代数和四元数环都属于这个范围；而像 $m\mathbb Z$（$m>1$）这样的环没有单位元，本节暂时不把它放进特征的定义中。

设 $R$ 是含幺环。定义映射
\[
\varphi:\mathbb Z\to R,\qquad n\mapsto n1_R.
\]
这里 $n1_R$ 表示 $1_R$ 的 $n$ 倍：正整数时是反复相加，$0$ 时为 $0_R$，负整数时取加法逆元。
这个映射保持加法，因为
\[
\varphi(m+n)=(m+n)1_R=m1_R+n1_R=\varphi(m)+\varphi(n).
\]
它也保持乘法。直观地说，$(m1_R)(n1_R)$ 就是把 $m1_R$ 重复加 $n$ 次，结果为 $(mn)1_R$；因此
\[
\varphi(mn)=(mn)1_R=(m1_R)(n1_R)=\varphi(m)\varphi(n).
\]
并且 $\varphi(1)=1_R$。
所以 $\varphi$ 是一个保持单位元的环同态。

\begin{definition}[含幺环的特征]\label{def:ring-characteristic}
设 $R$ 是含幺环，并令
\[
\varphi:\mathbb Z\to R,\qquad n\mapsto n1_R.
\]
由于 $\ker\varphi$ 是 $\mathbb Z$ 的加法子群，必形如 $m\mathbb Z$，其中 $m\geq0$。
称这个整数 $m$ 为 $R$ 的\textbf{特征}，记为
\[
\operatorname{char}(R)=m.
\]
\end{definition}

这个定义可以换成更直观的说法。
若存在正整数 $m$ 使得
\[
m1_R=0_R,
\]
则 $\operatorname{char}(R)$ 就是满足这个等式的最小正整数。
若不存在这样的正整数，则 $\ker\varphi=\{0\}$，也就是 $\operatorname{char}(R)=0$。
因此环的特征仍然是在测量"$1_R$ 加多少次第一次回到 $0_R$"；只是现在对象从域扩大到了含幺环。

几个基本例子可以帮助区分域和一般环的情况。
在 $\mathbb Z$ 中，$1$ 加任意正整数次都不等于 $0$，所以
\[
\operatorname{char}(\mathbb Z)=0.
\]
在剩余类环 $\mathbb Z/n\mathbb Z$ 中，$n$ 个 $\overline{1}$ 相加等于 $\overline{0}$，而更小的正整数不会做到这一点，所以
\[
\operatorname{char}(\mathbb Z/n\mathbb Z)=n.
\]
特别地，$\mathbb Z/6\mathbb Z$ 的特征是 $6$。
这和域的情形不同：一般环的特征不一定是素数。
原因也很清楚，$\mathbb Z/6\mathbb Z$ 中有零因子，例如 $\overline2\cdot\overline3=\overline0$。

矩阵环的特征则由系数环决定。
更一般地，若 $R$ 是含幺环，则 $M_n(R)$ 的单位元是单位矩阵 $I_n$，其中对角线上每一项都是 $1_R$。
于是
\[
mI_n=0
\]
当且仅当这个零矩阵的每个对角元都为 $0_R$，也就是
\[
m1_R=0_R.
\]
所以
\[
\operatorname{char}(M_n(R))=\operatorname{char}(R).
\]
特别地，取 $R=F$ 为域，就得到 $\operatorname{char}(M_n(F))=\operatorname{char}(F)$。

同理，若 $G$ 是有限群，则群代数 $F[G]$ 的单位元是 $1_F e$，其中 $e$ 是群 $G$ 的单位元。
于是
\[
m(1_F e)=(m1_F)e.
\]
由于群代数以 $G$ 中元素为一组 $F$-基，基向量 $e$ 前的系数为零当且仅当 $m1_F=0$。
因此
\[
\operatorname{char}(F[G])=\operatorname{char}(F).
\]

从同态的角度看，$\varphi:\mathbb Z\to R$ 的核已经带有理想的雏形。
这里暂时只需要知道：核记录了哪些整数倍的 $1_R$ 变成 $0_R$，而 $\mathbb Z$ 的加法子群形如 $m\mathbb Z$，所以这个记录最终可以压缩成一个整数 $m$。

\begin{theorem}[整环的特征]\label{thm:domain-characteristic}
设 $R$ 是整环。
若存在 $a\in R$，$a\neq0$，以及正整数 $m$，使得
\[
ma=0,
\]
则存在素数 $p$，使得对任意 $b\in R$ 都有
\[
pb=0.
\]
特别地，整环的特征只能是 $0$ 或素数。
\end{theorem}

\begin{proof}
先证明一个更强的中间结论：由某个非零元 $a$ 满足 $ma=0$，可以推出所有元素都满足 $mb=0$。
任取 $b\in R$。
若 $b=0$，结论显然成立。
若 $b\neq0$，则
\[
0=(ma)b=a(mb).
\]
这里第一等号来自 $ma=0$，第二个等号来自 $ma$ 的定义：$ma$ 是 $a$ 的 $m$ 倍，$(ma)b$ 就是把 $b$ 乘上这个倍数。
因为 $R$ 是整环，没有零因子；又 $a\neq0$，所以必须有
\[
mb=0.
\]
于是对任意 $b\in R$ 都有 $mb=0$。

特别地，取 $b=1_R$，得到
\[
m1_R=0.
\]
所以 $R$ 有正特征。设 $p$ 是使得 $p1_R=0$ 的最小正整数。
下面证明 $p$ 必为素数。

反设 $p$ 不是素数，则可以写成
\[
p=qr,\qquad q,r>1.
\]
由 $p1_R=0$ 得
\[
0=p1_R=(qr)1_R=(q1_R)(r1_R).
\]
整环无零因子，所以
\[
q1_R=0\qquad\text{或}\qquad r1_R=0.
\]
但 $q,r$ 都是小于 $p$ 的正整数，这与 $p$ 的最小性矛盾。
因此 $p$ 是素数。

最后，对任意 $b\in R$，
\[
pb=(p1_R)b=0\cdot b=0.
\]
这说明一旦整环中有非零元素在加法下具有有限阶，整个整环都会落在同一个素特征 $p$ 下。
反过来说，如果没有任何正整数使 $1_R$ 归零，那么整环的特征就是 $0$。
所以整环的特征只能是 $0$ 或素数。
\end{proof}

\section{模与代数}\label{sec:modules-and-algebras}

\subsection{模的定义}\label{sec:module-definition}

模可以看成线性空间的推广。
在线性空间里，向量可以相加，标量可以相加和相乘，标量还可以作用在向量上。
如果把标量所在的域换成一般的环，就得到模的概念。

这里涉及两个集合和一个外部作用。
被作用的集合通常记为 $M$，它至少要能做加法，并且在加法下构成阿贝尔群。
作用集合是环 $R$，它自身有加法和乘法。
外部作用是一个映射
\[
R\times M\to M,\qquad (r,x)\mapsto rx,
\]
这就是模里最容易混淆的地方：一个模实际牵涉到至少四种运算，分别是 $M$ 上的加法、$R$ 上的加法和乘法、以及 $R$ 对 $M$ 的作用。

\begin{definition}[左模]\label{def:left-module}
设 $R$ 是一个环。
一个\textbf{左 $R$-模}是一个阿贝尔群 $(M,+)$，配备一个作用
\[
R\times M\to M,\qquad (r,x)\mapsto rx,
\]
满足对任意 $r,s\in R$ 和 $x,y\in M$：
\begin{enumerate}
    \item $r(x+y)=rx+ry$；
    \item $(r+s)x=rx+sx$；
    \item $(rs)x=r(sx)$。
\end{enumerate}
如果 $R$ 有单位元 $1_R$，并且还满足
\[
1_Rx=x,\qquad x\in M,
\]
则称它是一个含幺左 $R$-模。
\end{definition}

前两条是两种分配律：环元素作用在 $M$ 的加法上要能拆开，环自身加法作用在同一个元素上也要能拆开。
第三条是乘法相容性。
它说 $rs$ 先在环里乘完再作用到 $x$ 上，和先让 $s$ 作用到 $x$，再让 $r$ 作用，结果相同：
\[
(rs)x=r(sx).
\]
这条公理 $(rs)x = r(sx)$ 正是"左"的含义所在：环中先乘 $rs$ 再作用到 $x$，等于先让 $s$ 作用、再让 $r$ 作用——作用顺序与乘法顺序相反。

右模的定义类似，只是作用写在右边：
\[
M\times R\to M,\qquad (x,r)\mapsto xr.
\]
对应的乘法相容性变成
\[
(xr)s=x(rs).
\]
也就是说，左模的乘法相容性 $(rs)x = r(sx)$ 表现为：先让 $s$ 作用到 $x$，再让 $r$ 作用——与环中 $rs$ 的书写顺序相反。
右模的乘法相容性 $(xr)s = x(rs)$ 则表现为：先让 $r$ 作用，再让 $s$ 作用——与环中 $rs$ 的书写顺序一致。
如果 $R$ 是交换环，左作用和右作用通常可以互相转换，因此两者差别不明显。
但在非交换环上，这个区别不能省略。

双模也是沿着这个方向继续推广：同一个阿贝尔群可以同时有左侧环作用和右侧环作用，并要求两边作用彼此相容。
本讲义暂时不展开双模。
在几何与分析中，模的语言也会自然出现：例如开集上的光滑函数环可以被微分算子环作用，前者就成为后者的模。

\subsection{模的例子}\label{sec:module-examples}

\paragraph*{例 8：阿贝尔群就是 $\mathbb Z$-模}
任意阿贝尔群 $A$ 都可以看成一个 $\mathbb Z$-模。
整数 $n$ 对元素 $a\in A$ 的作用定义为倍加：
\[
na=\underbrace{a+\cdots+a}_{n\text{ 次}},\qquad n>0,
\]
\[
0a=0,\qquad (-n)a=-(na).
\]
这些定义和整数加法、乘法相容，所以给出了一个 $\mathbb Z$-模结构。
具体说，$n(a+b)=na+nb$ 来自阿贝尔群中重复加法可以逐项拆开；
\[
(m+n)a=ma+na
\]
对应整数加法；
\[
(mn)a=m(na)
\]
对应整数乘法。
当 $m,n$ 中有负数时，利用 $(-n)a=-(na)$ 的定义直接验证即可。
反过来，任意 $\mathbb Z$-模只看它自身的加法，就得到一个阿贝尔群。
因此，$\mathbb Z$-模和阿贝尔群本质上是同一类对象。

\paragraph*{例 9：多项式环通过线性变换作用}
设 $M$ 是实数域上的线性空间，$T:M\to M$ 是线性映射。
令
\[
R=\mathbb R[t]
\]
为实系数一元多项式环。
对多项式
\[
f(t)=a_0+a_1t+\cdots+a_nt^n
\]
定义它对 $x\in M$ 的作用为
\[
f(t)\cdot x=f(T)(x)
=a_0x+a_1T(x)+\cdots+a_nT^n(x).
\]
这样，$M$ 就成为一个 $\mathbb R[t]$-模。
先看两条分配律。
由于 $f(T)$ 是线性映射，
\[
f(t)\cdot(x+y)=f(T)(x+y)=f(T)(x)+f(T)(y).
\]
而
\[
(f+g)(t)\cdot x=(f+g)(T)(x)=f(T)(x)+g(T)(x).
\]
乘法相容性则来自多项式代入线性变换的基本规则：
\[
(fg)(T)=f(T)g(T).
\]
因此
\[
(fg)(t)\cdot x=(fg)(T)(x)=f(T)(g(T)(x))=f(t)\cdot(g(t)\cdot x).
\]
常数多项式 $1$ 代入 $T$ 后是恒等映射，所以 $1\cdot x=x$。

这个例子很重要，因为它把"一个线性空间加上一个线性变换"翻译成了"一个多项式环上的模"。
例如，$T$ 的不变子空间恰好就是 $R$-子模，特征多项式和最小多项式的研究也可以在模的框架下统一处理。

\paragraph*{例 10：矩阵环作用在列向量上}
设 $A$ 是交换环。
矩阵环 $M_n(A)$ 可以自然作用在 $A^n$ 上：
\[
B\cdot x=Bx,\qquad B\in M_n(A),\ x\in A^n,
\]
右边就是通常的矩阵乘列向量。
模公理逐条对应矩阵乘法的基本性质：
\[
B(x+y)=Bx+By,
\]
\[
(B+C)x=Bx+Cx,
\]
而矩阵乘法结合律正好保证
\[
(BC)x=B(Cx),
\]
若使用含幺模约定，单位矩阵 $I_n$ 还满足 $I_nx=x$。
因此 $A^n$ 是一个左 $M_n(A)$-模。

\paragraph*{例 11：表示论中的模}
群表示也可以用模的语言重写。
设 $G$ 是有限群，$F$ 是域。
一个 $F$ 上的线性表示，就是给出一个 $F$-向量空间 $V$，以及一个群同态
\[
\rho:G\to \operatorname{GL}(V).
\]
也就是说，每个群元素 $g$ 都对应一个可逆线性变换 $\rho(g)$，并且满足
\[
\rho(gh)=\rho(g)\rho(h)
\]
和 $\rho(e)=\operatorname{id}_V$。

表示一开始只告诉我们：每个群元素 $g$ 怎么作用在 $V$ 上。
群代数 $F[G]$ 的元素是群元素的线性组合，所以最自然的想法是：让每个群元素先按表示作用，再把结果按系数线性组合。

对群代数中的元素
\[
\alpha=\sum_{g\in G}a_g g
\]
和向量 $v\in V$，定义
\[
\alpha\cdot v=\sum_{g\in G}a_g\rho(g)(v).
\]
这就是把基元素的作用
\[
g\cdot v=\rho(g)(v)
\]
按线性性扩展到整个群代数。

先看加法方面。
因为 $V$ 是 $F$-向量空间，线性组合可以逐项拆开，所以
\[
\alpha\cdot(v+w)=\alpha\cdot v+\alpha\cdot w
\]
和
\[
(\alpha+\beta)\cdot v=\alpha\cdot v+\beta\cdot v
\]
都直接来自 $\rho(g)$ 的线性性和系数加法。

真正需要看的是乘法相容性。
先看基元素。若 $g,h\in G$，则
\[
(gh)\cdot v=\rho(gh)(v)=\rho(g)(\rho(h)(v))=g\cdot(h\cdot v).
\]
这里用到的正是表示的同态条件 $\rho(gh)=\rho(g)\rho(h)$。

一般元素只是把这个等式按双线性展开。
若
\[
\alpha=\sum_{g\in G}a_g g,\qquad \beta=\sum_{h\in G}b_h h,
\]
那么 $\alpha\beta$ 中由一对 $(g,h)$ 贡献的项是 $a_gb_h(gh)$。
这个基元素作用到 $v$ 上，给出
\[
a_gb_h\rho(gh)(v)=a_gb_h\rho(g)(\rho(h)(v)).
\]
把所有 $g,h$ 的贡献加起来，就得到
\[
(\alpha\beta)\cdot v
=\sum_{g,h\in G}a_gb_h\rho(gh)(v)
=\sum_{g,h\in G}a_gb_h\rho(g)(\rho(h)(v))
=\alpha\cdot(\beta\cdot v).
\]
最后一个等号来自：$\beta\cdot v=\sum_h b_h\rho(h)(v)$，再让 $\alpha$ 作用即得。

群代数的单位元是 $1_F e$，而 $\rho(e)=\operatorname{id}_V$，所以
\[
(1_F e)\cdot v=v.
\]
因此表示空间 $V$ 成为一个左 $F[G]$-模。
证明的核心是群代数乘法中的 $gh$ 与表示条件 $\rho(gh)=\rho(g)\rho(h)$ 的吻合。

反过来，若给定一个左 $F[G]$-模 $V$，把群元素 $g\in G$ 看作 $F[G]$ 中的基元素，它对 $V$ 的作用就是一个线性变换。
由于 $g^{-1}$ 也在群中，$g$ 的作用有逆作用，所以得到 $\rho(g)\in\operatorname{GL}(V)$。
模公理中的乘法相容性给出
\[
\rho(gh)=\rho(g)\rho(h).
\]
所以 $F[G]$-模和 $G$ 的线性表示本质上是在说同一件事。
这把"群元素如何线性地作用"统一放进了环作用的框架里。

类似地，quiver 表示也和路径代数上的模密切相关。
在每个顶点放一个向量空间、每条箭头放一个线性映射，这些数据可以由路径代数的作用组织起来。
这个联系会在表示论中继续展开。

\subsection{代数的定义}\label{sec:algebra-definition}

在此基础上，再引入代数的概念。
如果 $A$ 只是一个 $R$-模，那么 $A$ 中的元素可以相加，也可以被 $R$ 中元素作用，但 $A$ 自己的元素之间未必能相乘。
如果 $A$ 自身还有一个乘法，并且这个乘法和 $R$-模结构相容，就得到 $R$-代数。

\begin{definition}[$R$-代数]\label{def:algebra}
设 $R$ 是交换含幺环。
一个\textbf{$R$-代数}是一个 $R$-模 $A$，同时配备一个乘法
\[
A\times A\to A,\qquad (x,y)\mapsto xy,
\]
使得 $A$ 在这个乘法下成为环，并且乘法对 $R$-模结构双线性：
\[
r(x+y)=rx+ry,\qquad r(xy)=(rx)y=x(ry),
\]
其中 $r\in R$，$x,y\in A$。
若 $A$ 有乘法单位元，还通常要求 $R$ 的单位元作用为 $A$ 的单位元。
\end{definition}

概括为一句话：
\[
\text{$R$-代数}=\text{带有相容乘法的 $R$-模}.
\]
它比单纯的 $R$-模多了 $A$ 内部的乘法，也比单纯的环多了来自 $R$ 的标量作用。

\paragraph*{例 12：复数作为实代数}
复数集 $\mathbb C$ 首先是一个 $\mathbb R$-向量空间，也就是一个 $\mathbb R$-模。
同时，复数之间本来就可以相乘。
作为环，$\mathbb C$ 的结合律、分配律和单位元都来自通常的复数乘法。
作为 $\mathbb R$-模，它的标量作用就是实数乘复数。
最后要检查的是这两套结构是否相容：
\[
r(zw)=(rz)w=z(rw),\qquad r\in\mathbb R,\ z,w\in\mathbb C.
\]
这个等式成立，是因为实数 $r$ 在复数乘法中可以看成复数 $r+0i$，并且它与任意复数相乘都交换、结合。
因此 $\mathbb C$ 是一个 $\mathbb R$-代数。
这个例子说明，代数并不是又凭空多造了一个对象；它只是把"标量作用"和"对象内部的乘法"同时记在同一个结构里。

\paragraph*{补充：$K$-代数和 $R$-代数的区别}

有时书上写 $K$-代数，有时写 $R$-代数，二者的本质差别在于底层标量对象不同。
通常 $K$ 表示一个域，而 $R$ 表示一个交换含幺环。

如果 $K$ 是域，那么一个 $K$-代数 $A$ 首先是一个 $K$-向量空间。
也就是说，$A$ 中的元素可以按 $K$ 中的标量做线性组合。
同时 $A$ 自己还带有乘法，并且这个乘法要和 $K$-向量空间结构相容：
\[
\lambda(xy)=(\lambda x)y=x(\lambda y),\qquad
\lambda\in K,\ x,y\in A.
\]
所以可以把它记成
\[
\text{$K$-代数}=\text{带有相容乘法的 $K$-向量空间}.
\]

如果 $R$ 只是一般交换含幺环，那么一个 $R$-代数 $A$ 首先只是一个 $R$-模，而不一定是向量空间。
这时仍然要求 $A$ 有自己的乘法，并且满足
\[
r(xy)=(rx)y=x(ry),\qquad r\in R,\ x,y\in A.
\]
因此
\[
\text{$R$-代数}=\text{带有相容乘法的 $R$-模}.
\]

这说明 $K$-代数其实是 $R$-代数的一个特殊情形：当底环 $R$ 恰好是域 $K$ 时，$R$-模就是 $K$-向量空间。
域上的向量空间性质更好，例如基和维数理论完备；一般环上的模则可能有挠元，也可能不是自由模，所以 $R$-代数比 $K$-代数更一般，也更复杂。

从含幺代数的角度看，还可以用结构映射来理解。
给 $A$ 一个含幺 $R$-代数结构，等价于给一个环同态
\[
R\to Z(A),
\]
把 $R$ 中的标量送进 $A$ 的中心。
这样 $r\in R$ 才能在 $A$ 中像标量一样左右移动，保证
\[
(rx)y=x(ry).
\]
若把 $R$ 换成域 $K$，同样得到 $K\to Z(A)$，只是此时 $A$ 还自动带有 $K$-向量空间结构。

例如 $\mathbb C$ 是一个 $\mathbb R$-代数，因为复数可以看成实向量空间，并且复数乘法与实数标量乘法相容。
而 $\mathbb Z[x]$ 是一个 $\mathbb Z$-代数：它是一个环，也是一个 $\mathbb Z$-模，但这里的底层结构不是向量空间，而是整数模。
所以看到"$K$-代数"时，可以先想到"域上的线性结构加乘法"；看到"$R$-代数"时，则要提醒自己：标量可能只是一个环，底层对象只是模。

\section{同态、子结构与理想}\label{sec:hom-substructure-ideals}

\subsection{环和模的同态}\label{sec:ring-module-homomorphism}

有了环和模以后，下一步自然要问：两个这样的对象之间，什么样的映射才算"保持结构"？
在域论里，同态要保持加法、乘法和单位元；环和模的情形更灵活一些，因为环未必有单位元，模还带有外部作用。

\begin{definition}[环同态]\label{def:ring-homomorphism}
设 $R,S$ 是两个环。
映射
\[
f:R\to S
\]
称为一个\textbf{环同态}，如果对任意 $a,b\in R$ 都有
\[
f(a+b)=f(a)+f(b),\qquad f(ab)=f(a)f(b).
\]
若 $R,S$ 都有单位元，并且还满足
\[
f(1_R)=1_S,
\]
则称 $f$ 是\textbf{保持单位元的环同态}。
\end{definition}

关于单位元的条件，不同教材有不同约定。
有的教材把 $f(1_R)=1_S$ 放进环同态的定义里；也有教材只要求保持加法和乘法，把保持单位元作为额外条件。
本讲义在需要单位元参与的地方会明确说明"保持单位元"。
这样做的好处是可以同时容纳无单位元环，也能解释一些自然出现但不保持单位元的映射。

环同态会自动保持零元。
因为
\[
f(0_R)=f(0_R+0_R)=f(0_R)+f(0_R),
\]
两边同时加上 $-f(0_R)$，得到 $f(0_R)=0_S$。
同样地，环同态也保持加法逆元：
\[
0_S=f(0_R)=f(a+(-a))=f(a)+f(-a),
\]
所以 $f(-a)=-f(a)$。

\begin{definition}[模同态]\label{def:module-homomorphism}
设 $M,N$ 是同一个环 $R$ 上的左模。
映射
\[
T:M\to N
\]
称为一个\textbf{$R$-模同态}，如果对任意 $x,y\in M$ 和 $r\in R$ 都有
\[
T(x+y)=T(x)+T(y),\qquad T(rx)=rT(x).
\]
\end{definition}

第一条说明 $T$ 保持模自身的加法；第二条说明 $T$ 和 $R$ 的作用相容。
如果把线性空间看成域上的模，那么模同态就是线性映射。
所以模同态并不是新造的概念，而是把"线性映射"从域上的向量空间推广到环上的模。

\paragraph*{环同态的例子}
一个典型的环同态来自剩余类环。
设 $m,n$ 是正整数，定义
\[
\pi:\mathbb Z/(mn)\mathbb Z\to \mathbb Z/m\mathbb Z,\qquad
[a]_{mn}\mapsto [a]_m.
\]
这个映射是良定义的：如果 $a\equiv b\pmod{mn}$，那么 $mn\mid(a-b)$，特别地 $m\mid(a-b)$，所以 $a\equiv b\pmod m$。
它显然保持加法和乘法，也保持单位元。
因此 $\pi$ 是一个满的环同态。

反方向一般没有这样的自然同态。
如果想把
\[
\mathbb Z/m\mathbb Z\to \mathbb Z/(mn)\mathbb Z
\]
也写成 $[a]_m\mapsto[a]_{mn}$，就会遇到良定义问题：在 $\mathbb Z/m\mathbb Z$ 中，$[0]_m=[m]_m$，但在 $\mathbb Z/(mn)\mathbb Z$ 中，$[0]_{mn}$ 和 $[m]_{mn}$ 通常不同。
若要求保持单位元，这个障碍更直接：从 $\mathbb Z/m\mathbb Z$ 出发时，$m\cdot 1=0$，但在 $\mathbb Z/(mn)\mathbb Z$ 中通常有 $m\cdot1\neq0$。

\paragraph*{模同态的例子}
模同态的例子可以从更熟悉的 $\mathbb Z$-模来看。
把阿贝尔群都看成 $\mathbb Z$-模，则自然投影
\[
q:\mathbb Z\to \mathbb Z/m\mathbb Z,\qquad a\mapsto [a]_m
\]
是一个 $\mathbb Z$-模同态。
它保持加法，并且整数倍作用与映射相容：
\[
q(ra)=[ra]_m=r[a]_m=rq(a).
\]
这个例子中
\[
\ker q=m\mathbb Z,\qquad \operatorname{im}q=\mathbb Z/m\mathbb Z.
\]

和群、域的情形一样，环同态和模同态也可以分为单同态、满同态和同构。
对同态 $f:R\to S$，定义
\[
\ker f=\{a\in R:f(a)=0\},\qquad
\operatorname{im}f=\{f(a):a\in R\}.
\]
同态是单射，当且仅当 $\ker f=\{0\}$；同态是满射，当且仅当 $\operatorname{im}f=S$。
对模同态 $T:M\to N$，核与像定义为
\[
\ker T=\{x\in M:T(x)=0_N\},\qquad
\operatorname{im}T=\{T(x):x\in M\}.
\]
其中 $\ker T$ 是 $M$ 的子模，$\operatorname{im}T$ 是 $N$ 的子模。
模同态 $T$ 是单射，当且仅当 $\ker T=\{0\}$；是满射，当且仅当 $\operatorname{im}T=N$。

\subsection{子结构：子环与子模}\label{sec:substructures}

同态告诉我们怎样在两个对象之间传递结构；子结构则回答另一个问题：一个对象内部的哪些子集本身仍然是同类对象？

\begin{definition}[子环]\label{def:subring}
设 $R$ 是环，$S\subseteq R$。
若 $S$ 在 $R$ 的加法和乘法下本身构成一个环，则称 $S$ 是 $R$ 的\textbf{子环}。
\end{definition}

本讲义采用的约定是：子环不要求包含原环的乘法单位元。
这和理想的讨论更相容，因为理想通常不会含有 $1$；一旦某个理想含有 $1$，它就会等于整个环。
例如 $2\mathbb Z$ 是 $\mathbb Z$ 的子环，但它不含 $\mathbb Z$ 的单位元 $1$。

\begin{theorem}[子环判定]\label{thm:subring-test}
设 $S\subseteq R$，若
\[
S\neq\varnothing,\qquad a-b\in S,\qquad ab\in S\quad(a,b\in S),
\]
则 $S$ 是 $R$ 的子环。
\end{theorem}

\begin{proof}
因为 $S\neq\varnothing$，可取 $a\in S$。
由 $a-b\in S$，取 $b=a$ 得
\[
0=a-a\in S.
\]
再取 $a=0$，对任意 $b\in S$，得到
\[
-b=0-b\in S.
\]
于是对任意 $a,b\in S$，
\[
a+b=a-(-b)\in S.
\]
所以 $S$ 对加法封闭，并且含有零元和每个元素的加法逆元。
条件 $ab\in S$ 保证乘法封闭。
结合律和分配律不需要重新证明，因为它们直接继承自 $R$。
因此 $S$ 是 $R$ 的子环。
\end{proof}

\begin{definition}[子模]\label{def:submodule}
设 $M$ 是左 $R$-模，$N\subseteq M$。
若 $N$ 在 $M$ 的加法下构成子阿贝尔群，并且对任意 $r\in R$、$x\in N$ 都有
\[
rx\in N,
\]
则称 $N$ 是 $M$ 的\textbf{子模}。
\end{definition}

子模的条件比子环更像线性代数中的子空间：不仅要对加法和取负封闭，还要对标量作用封闭。
若 $R$ 是域，这正是子空间的定义。
若 $R$ 是一般环，子模则是子空间概念的模论版本。

现在把环 $R$ 看成它自身上的左模，作用就是左乘。
这时一个子模 $I\subseteq R$ 要满足
\[
a-b\in I,\qquad r a\in I\quad(r\in R,\ a\in I).
\]
这已经非常接近理想的定义。
若同时要求右乘也封闭，即 $ar\in I$，就得到后面要用的双边理想。
所以理想可以理解为：当环 $R$ 看作自身上的左模时，那些同时对左乘和右乘封闭的子模。

前面已经讲过群表示和 $F[G]$-模的关系。
从子结构角度看，表示中的不变子空间正是对应模的子模：一个子空间如果被所有群元素作用后仍留在自身中，就等价于它对整个群代数 $F[G]$ 的作用封闭。
这正是模语言把表示论组织起来的原因。

\subsection{核、像与理想}\label{sec:kernels-images-ideals}

同态的核与像都来自同一个问题：一个结构经过映射以后，哪些东西被送到 $0$，哪些东西真的出现在目标里？
但在环论里，核比像更特别。
它不只是一个子环，还会被整个定义域的乘法吸收。
这个额外性质正是理想的来源。

设 $f:R\to S$ 是环同态。
先看核
\[
\ker f=\{a\in R:f(a)=0\}.
\]
它是 $R$ 的子环。若 $a,b\in\ker f$，则
\[
f(a-b)=f(a)-f(b)=0-0=0,
\]
所以 $a-b\in\ker f$；又
\[
f(ab)=f(a)f(b)=0\cdot0=0,
\]
所以 $ab\in\ker f$。

像
\[
\operatorname{im}f=\{f(a):a\in R\}
\]
也是 $S$ 的子环。
例如两个像元素 $f(a),f(b)$ 的差和积分别是
\[
f(a)-f(b)=f(a-b),\qquad f(a)f(b)=f(ab),
\]
仍然落在像中。

现在看核真正多出来的性质。
若 $a\in\ker f$，$r\in R$，则
\[
f(ra)=f(r)f(a)=f(r)\cdot0=0,
\]
所以 $ra\in\ker f$。
同理，
\[
f(ar)=f(a)f(r)=0\cdot f(r)=0,
\]
所以 $ar\in\ker f$。
也就是说，核不仅对自身的加减乘封闭，还能被整个环从左、右两边相乘而不跑出去。
既然这个性质会反复出现，我们就把满足这种吸收性的加法子群单独命名。

\begin{definition}[理想]\label{def:ideal}
设 $R$ 是环，$I\subseteq R$。
若 $I$ 是 $R$ 的加法子群，并且对任意 $r\in R$、$a\in I$ 都有
\[
ra\in I,\qquad ar\in I,
\]
则称 $I$ 是 $R$ 的\textbf{双边理想}，简称\textbf{理想}。
\end{definition}

如果 $R$ 是交换环，左右乘法没有区别，所以只需写
\[
rI\subseteq I,\qquad r\in R.
\]
若 $R$ 非交换，则需要区分左理想、右理想和双边理想。
左理想只要求 $ra\in I$，右理想只要求 $ar\in I$；能用来构造商环的通常是双边理想。

上面的计算说明：环同态的核一定是理想。
这也解释了为什么前面讨论环的特征时，映射
\[
\varphi:\mathbb Z\to R,\qquad n\mapsto n1_R
\]
的核会自然出现。
它不仅是 $\mathbb Z$ 的一个加法子群，更是 $\mathbb Z$ 的一个理想；而 $\mathbb Z$ 的理想都形如 $m\mathbb Z$，于是特征就被一个整数记录下来。

\subsection{理想的基本性质与主理想}\label{sec:ideal-basic-properties}

理想的定义看起来只是一个封闭性条件，但它对环的结构限制很强。
最简单的性质是：一旦理想里出现单位元，它就不再是一个"小"的子结构。

\begin{proposition}[含有单位元的理想]\label{prop:ideal-containing-one}
设 $R$ 是含幺环，$I$ 是 $R$ 的理想。
若 $1_R\in I$，则 $I=R$。
\end{proposition}

\begin{proof}
任取 $r\in R$。
因为 $1_R\in I$，且理想对环中元素的乘法吸收，所以
\[
r=r1_R\in I.
\]
因此 $R\subseteq I$。
反过来 $I\subseteq R$ 本来就成立，所以 $I=R$。
\end{proof}

这个性质在域里给出一个非常强的结论。

\begin{proposition}[域只有平凡理想]\label{prop:field-ideals}
设 $F$ 是域。
则 $F$ 的理想只有 $\{0\}$ 和 $F$。
\end{proposition}

\begin{proof}
设 $I$ 是 $F$ 的理想。
若 $I=\{0\}$，已经是平凡理想。
否则，取非零元素 $a\in I$。
因为 $F$ 是域，$a$ 有乘法逆元 $a^{-1}$。
由理想的吸收性，
\[
1=a^{-1}a\in I.
\]
于是由命题 \ref{prop:ideal-containing-one} 得 $I=F$。
\end{proof}

这和域同态必为单射的现象是同一件事的两个侧面。
域同态的核是理想，而域的理想只有 $\{0\}$ 和整个域；如果同态保持单位元，核不可能是整个域，所以只能是 $\{0\}$。

在交换环中，最常见的一类理想由单个元素生成。
设 $R$ 是交换环，$a\in R$，定义
\[
(a)=aR=\{ar:r\in R\}.
\]
这是一个理想。
若 $ar,as\in aR$，则
\[
ar-as=a(r-s)\in aR;
\]
对任意 $t\in R$，
\[
t(ar)=a(tr)\in aR.
\]
所以 $aR$ 对加减和环中元素乘法都封闭。
这样的理想称为由 $a$ 生成的\textbf{主理想}。

"理想"这个名字来自代数数论中对唯一分解失败的修正。
早期说的是"理想数"，后来对象本身逐渐抽象成环中的特殊子集，名称也就简化为"理想"。
这里不展开历史，只需要记住它的代数功能：理想是环同态的核，也是构造商环时用来做"除法"的对象。

整数环中的理想是最基本的主理想例子。

\begin{theorem}[$\mathbb Z$ 的理想]\label{thm:z-ideals-principal}
$\mathbb Z$ 中每个理想都形如 $m\mathbb Z$。
\end{theorem}

\begin{proof}
设 $I$ 是 $\mathbb Z$ 的理想。
若 $I=\{0\}$，则 $I=0\mathbb Z$。
下面设 $I\neq\{0\}$。

因为 $I$ 是加法子群，若 $n\in I$，则 $-n\in I$。
所以 $I$ 中存在正整数。
取 $I$ 中最小的正整数 $m$。
由于 $m\in I$，理想对整数乘法封闭，所以
\[
m\mathbb Z\subseteq I.
\]

反过来，任取 $n\in I$。
用带余除法写
\[
n=qm+r,\qquad 0\leq r<m.
\]
因为 $n\in I$，$qm\in I$，所以
\[
r=n-qm\in I.
\]
若 $r>0$，就得到一个比 $m$ 更小的正整数属于 $I$，与 $m$ 的最小性矛盾。
因此 $r=0$，即 $m\mid n$。
所以 $n\in m\mathbb Z$。

综上，$I=m\mathbb Z$。
\end{proof}

\paragraph*{例：一元多项式环}

$\mathbb Z$ 中每个理想都是主理想，这来自整数的带余除法。
同样的思想也出现在一元多项式环中。
这里以 $\mathbb C[x]$ 为例；事实上，任意域 $F$ 上的一元多项式环 $F[x]$ 都有相同性质。

\begin{theorem}[$\mathbb C[x]$ 是主理想整环]\label{thm:cx-pid}
$\mathbb C[x]$ 中每个理想都是主理想。
\end{theorem}

\begin{proof}
设 $I$ 是 $\mathbb C[x]$ 的理想。
若 $I=\{0\}$，则 $I=(0)$。
下面设 $I\neq\{0\}$。

如果 $I$ 中含有非零常数 $c$，那么 $c$ 在 $\mathbb C[x]$ 中可逆，且 $c^{-1}c=1$。
于是 $1\in I$，从而 $I=\mathbb C[x]=(1)$。
所以对真理想来说，它不能含有非零常数。

现在考虑非零真理想的情形。
在 $I$ 的非零多项式中，取次数最小的一个，记为 $P(x)$。
我们证明 $I=(P)$。
由于 $P\in I$，显然有
\[
(P)\subseteq I.
\]

反过来，任取 $f(x)\in I$。
用多项式带余除法除以 $P(x)$：
\[
f(x)=q(x)P(x)+r(x),\qquad r(x)=0\ \text{或}\ \deg r<\deg P.
\]
因为 $f(x)\in I$，$q(x)P(x)\in I$，所以
\[
r(x)=f(x)-q(x)P(x)\in I.
\]
若 $r(x)\neq0$，它就是 $I$ 中次数比 $P$ 更低的非零多项式，矛盾。
因此 $r(x)=0$，即 $P$ 整除 $f$。
所以 $f\in(P)$。

于是 $I\subseteq(P)$，从而 $I=(P)$。
\end{proof}

这个证明和 $\mathbb Z$ 的证明完全平行：都取一个"最小"对象，然后用带余除法把任意元素除过去，余数若非零就违反最小性。
整数里"最小"指的是绝对值最小的正整数；一元多项式里"最小"指的是次数最低。

但多元多项式环就不再这么简单。
例如 $\mathbb C[x,y]$ 是整环，却不是主理想环。
考虑理想
\[
(x,y)=\{xf(x,y)+yg(x,y):f,g\in\mathbb C[x,y]\}.
\]
它由 $x$ 和 $y$ 共同生成。
如果它是主理想，设
\[
(x,y)=(h).
\]
因为 $x,y\in(h)$，所以 $h$ 同时整除 $x$ 和 $y$。
在 $\mathbb C[x,y]$ 中，$x$ 和 $y$ 没有非常数公因子，因此 $h$ 只能是非零常数。
但若 $h$ 是非零常数，则 $(h)=\mathbb C[x,y]$，会推出 $1\in(x,y)$。
这不可能，因为 $xf(x,y)+yg(x,y)$ 的常数项总是 $0$。
所以 $(x,y)$ 不是主理想。

这个例子说明，"整环"和"主理想环"是两种不同性质。
$\mathbb C[x,y]$ 没有零因子，但它的理想已经比一元多项式环复杂得多。
后面讨论商环时，理想的这种复杂性会继续出现。

\subsection{生成理想与双边理想}\label{sec:generated-two-sided-ideals}

在具体计算中，我们经常不是先拿到一个完整的理想，而是先指定若干元素，然后问：包含这些元素的最小理想是什么？
这就是生成理想。

圆括号 $(a)$ 常用来表示由 $a$ 生成的理想。
这个记号和商环中的等价类记号有时会同时出现，所以阅读时要看上下文：在"生成"语境中，$(a)$ 是理想；在"模掉某个理想"的语境中，$a+I$ 或 $\overline a$ 才是等价类。

若 $R$ 是交换含幺环，最常用的形式非常简单：
\[
(a_1,\dots,a_m)=Ra_1+\cdots+Ra_m.
\]
也就是说，$(a_1,\dots,a_m)$ 中的元素都是有限和
\[
r_1a_1+\cdots+r_ma_m,\qquad r_i\in R.
\]
它是包含 $a_1,\dots,a_m$ 的最小理想。
例如在 $\mathbb C[x,y]$ 中，
\[
(x,y)=\{xf(x,y)+yg(x,y):f,g\in\mathbb C[x,y]\}.
\]

非交换环里要更小心，因为左乘和右乘不一定相同。
如果要生成双边理想，就必须允许生成元左右两边都乘上环中元素。
当 $R$ 是含幺环时，由一个元素 $a$ 生成的双边理想为
\[
(a)=RaR,
\]
也就是所有有限和
\[
\sum_{\ell=1}^N r_\ell a s_\ell,\qquad r_\ell,s_\ell\in R.
\]
更一般地，由集合 $A\subseteq R$ 生成的双边理想为
\[
(A)=RAR
=\left\{\sum_{\ell=1}^N r_\ell a_\ell s_\ell:
r_\ell,s_\ell\in R,\ a_\ell\in A,\ N\in\mathbb N\right\}.
\]

这里含幺条件很重要。
若 $R$ 有单位元，那么
\[
a=1_Ra1_R,
\]
所以生成元本身已经包含在 $RaR$ 里。
但若 $R$ 没有单位元，就不能用 $1_Ra1_R$ 把 $a$ 放进去。
这时由 $A$ 生成的双边理想需要额外加入生成元本身及其左右倍数：
\[
\mathbb ZA+RA+AR+RAR.
\]
其中
\[
\mathbb ZA
=\left\{\sum_{\ell=1}^N n_\ell a_\ell:
n_\ell\in\mathbb Z,\ a_\ell\in A\right\}.
\]
所以这里真正需要记住的不是表格，而是差别来源：有单位元时，$a=1a1$；无单位元时，生成元本身不会自动出现在 $RAR$ 里，所以必须额外补进去。

\section{商环和商模}\label{sec:quotient-rings-and-modules}

\subsection{从理想到商环}\label{sec:quotient-ring-construction}

前面讲理想时，我们已经看到：理想是环同态的核，也是构造商环时用来做"除法"的对象。
现在把第二个作用展开。
直观地说，理想可以规定哪些元素应该被看成零；一旦这些元素被压成零，原来的环就会折叠成一个新的环。

最熟悉的例子是整数的模 $n$ 运算。
在 $\mathbb Z/n\mathbb Z$ 中，两个整数如果相差 $n$ 的倍数，就被看成同一个剩余类。
现在把"相差 $n$ 的倍数"换成"相差一个理想里的元素"。
设 $R$ 是环，$I$ 是 $R$ 的双边理想。对 $a,b\in R$，定义
\[
a\sim b \quad\Longleftrightarrow\quad a-b\in I.
\]
这一步本质上只用到了 $I$ 是加法子群。
因为加法子群对取负和加法封闭，所以这个关系确实给出了 $R$ 的一个划分。
元素 $a$ 所在的等价类记为
\[
a+I=\{a+x:x\in I\}.
\]
所有等价类组成的集合记为 $R/I$。

接下来要把 $R$ 中的加法和乘法传到等价类上。最自然的定义是
\[
(a+I)+(b+I)=(a+b)+I,
\qquad
(a+I)(b+I)=ab+I.
\]
这里真正需要检查的是良定义：如果换一个代表元，结果会不会变。

先看加法。若 $a-a'\in I$，$b-b'\in I$，则
\[
(a+b)-(a'+b')=(a-a')+(b-b')\in I.
\]
所以 $(a+b)+I=(a'+b')+I$，加法不依赖代表元的选择。

乘法要多一步。若仍有 $a-a'\in I$，$b-b'\in I$，则
\[
ab-a'b'=(a-a')b+a'(b-b').
\]
第一项 $(a-a')b$ 属于 $I$，用到的是 $I$ 对右乘的吸收性；第二项 $a'(b-b')$ 属于 $I$，用到的是 $I$ 对左乘的吸收性。
因此 $ab-a'b'\in I$，从而 $ab+I=a'b'+I$。
这正是双边理想在商环构造中出现的原因：在非交换环里，只要求左理想或右理想，一般不能保证商环乘法良定义。

\hypertarget{q:quotient-ring-two-sided}{}
\medskip
\noindent\textbf{思考：商环的加法良定义只用到 $I$ 是加法子群，乘法良定义为什么必须要求 $I$ 是双边理想？}（\hyperlink{ans:quotient-ring-two-sided}{跳转到解答}。）
\medskip

\begin{definition}[商环]\label{def:quotient-ring}
设 $R$ 是环，$I$ 是 $R$ 的双边理想。
在 $R/I$ 上定义
\[
(a+I)+(b+I)=(a+b)+I,\qquad (a+I)(b+I)=ab+I.
\]
这样得到的环称为 $R$ 对 $I$ 的\textbf{商环}，记为 $R/I$。
\end{definition}

当 $R$ 是交换环时，左右乘法没有区别，"双边"条件自动满足；此时只需说 $I$ 是 $R$ 的理想。

商环的零元是 $0+I$，也就是 $I$ 这个等价类本身。
若 $R$ 有单位元 $1_R$，则 $1_R+I$ 是 $R/I$ 的单位元，因为
\[
(1_R+I)(a+I)=a+I,\qquad (a+I)(1_R+I)=a+I.
\]
若 $R$ 是交换环，则
\[
(a+I)(b+I)=ab+I=ba+I=(b+I)(a+I),
\]
所以 $R/I$ 也是交换环。

\subsection{自然映射}\label{sec:canonical-map-quotient-ring}

商环一旦构造出来，就自带一个最自然的映射：
\[
\pi:R\to R/I,\qquad \pi(a)=a+I.
\]
它把 $R$ 中的每个元素送到自己所在的等价类。
这个映射记录了商环的构造过程：原来在 $R$ 里不同的元素，只要相差一个 $I$ 中的元素，在 $R/I$ 中就被送到同一个类里。

映射 $\pi$ 是满的环同态。
满射性很直接：$R/I$ 中任意元素本来就是某个 $a+I$，而 $\pi(a)=a+I$。
它保持加法，因为
\[
\pi(a+b)=(a+b)+I=(a+I)+(b+I)=\pi(a)+\pi(b).
\]
它也保持乘法，因为
\[
\pi(ab)=ab+I=(a+I)(b+I)=\pi(a)\pi(b).
\]
所以 $\pi$ 是一个满环同态。

它的核正好是 $I$：
\[
\ker\pi=\{a\in R:\pi(a)=0+I\}
=\{a\in R:a+I=I\}=I.
\]
最后一个等号的意思是：$a+I=0+I$ 当且仅当 $a-0=a\in I$。

因此，商环可以理解为"把 $I$ 中所有元素压成零以后得到的环"。
自然映射 $\pi$ 就是这个压缩过程本身。
后面讲环同态第一基本定理时，会看到任意环同态都可以分解成这种形式：先把核压成零，再进入同态的像。

\subsection{两个基本例子}\label{sec:quotient-ring-examples}

\paragraph*{例 1：$\mathbb Q[x]/(x^2-2)$}
设
\[
I=(x^2-2)\subseteq \mathbb Q[x].
\]
商环 $\mathbb Q[x]/I$ 的元素是形如 $f(x)+I$ 的等价类。
对任意 $f(x)\in\mathbb Q[x]$，用 $x^2-2$ 作带余除法：
\[
f(x)=q(x)(x^2-2)+(ax+b),
\qquad a,b\in\mathbb Q.
\]
这里余式 $ax+b$ 由带余除法唯一确定。
于是
\[
f(x)+I=ax+b+I.
\]
也就是说，商环中的每个元素都可以由一次式代表，而且这个一次式代表是唯一的。
如果 $ax+b+I=a'x+b'+I$，则
\[
(a-a')x+(b-b')\in (x^2-2).
\]
左边次数小于 $2$，除非它是零多项式，否则不可能是 $x^2-2$ 的倍数。
因此 $a=a'$ 且 $b=b'$。

在这个商环中，
\[
x^2+I=2+I,
\]
因为 $x^2-2\in I$。
所以 $x+I$ 的行为像一个平方等于 $2$ 的元素。
这已经很像 $\sqrt2$ 在 $\mathbb Q(\sqrt2)$ 中的行为。

定义映射
\[
\overline\Phi:\mathbb Q[x]/(x^2-2)\to\mathbb Q(\sqrt2),
\qquad
\overline\Phi(f(x)+I)=f(\sqrt2).
\]
这个映射把 $x+I$ 送到 $\sqrt2$。
它是良定义的：若 $f(x)+I=g(x)+I$，则 $f(x)-g(x)\in I$，所以 $f(x)-g(x)$ 是 $x^2-2$ 的倍数；代入 $x=\sqrt2$ 后得到 $f(\sqrt2)-g(\sqrt2)=0$。

带余除法说明 $\mathbb Q[x]/I$ 中每个元素都可以写成 $ax+b+I$，而
\[
\overline\Phi(ax+b+I)=a\sqrt2+b.
\]
这正好给出 $\mathbb Q(\sqrt2)$ 中的元素。
反过来，$\mathbb Q(\sqrt2)$ 中每个 $b+a\sqrt2$ 都来自 $ax+b+I$。
因此
\[
\mathbb Q[x]/(x^2-2)\cong\mathbb Q(\sqrt2).
\]

这个例子说明，商环可以把"加入一个平方等于 $2$ 的元素"这件事代数化。
我们不是直接把 $\sqrt2$ 塞进 $\mathbb Q$，而是在多项式环里规定 $x^2-2$ 等于零。

\paragraph*{例 2：$\mathbb R[x]/(x^2+1)$}
复数也可以用同样的方法得到。
设
\[
I=(x^2+1)\subseteq \mathbb R[x].
\]
对任意 $g(x)\in\mathbb R[x]$，用 $x^2+1$ 作带余除法：
\[
g(x)=(x^2+1)p(x)+(ax+b),
\qquad a,b\in\mathbb R.
\]
所以每个等价类都可以写成
\[
ax+b+I.
\]

这类代表元的加法很直接：
\[
(ax+b+I)+(cx+d+I)=(a+c)x+(b+d)+I.
\]
乘法则需要用到商环中的关系 $x^2+I=-1+I$。
先在多项式环里相乘：
\[
(ax+b)(cx+d)=acx^2+(ad+bc)x+bd.
\]
因为 $x^2+1\in I$，所以 $x^2+I=-1+I$。
于是
\[
(ax+b+I)(cx+d+I)=(ad+bc)x+(bd-ac)+I.
\]

现在把
\[
ax+b+I
\]
对应到复数
\[
b+ai.
\]
加法也对应得上：
\[
(b+ai)+(d+ci)=(b+d)+(a+c)i,
\]
这与商环里的
\[
(ax+b+I)+(cx+d+I)=(a+c)x+(b+d)+I
\]
一致。
那么两个复数相乘得到
\[
(b+ai)(d+ci)=(bd-ac)+(ad+bc)i.
\]
这与上面的商环乘法完全一致。
所以这个对应同时保持加法和乘法。
因此
\[
\mathbb R[x]/(x^2+1)\cong\mathbb C.
\]

这个同构给出了复数域的一种等价定义。
复数不是凭空多出来的对象；它也可以看成在实系数多项式环中强行加入关系 $x^2+1=0$ 后得到的商环。
这正是商环构造的基本作用：把某个代数关系压进原来的环里。

\subsection{商模}\label{sec:quotient-modules}

商环是按理想取商；模也可以按子模取商。
这个构造和线性代数里的商空间非常接近：把一个子空间里的向量看成零向量，然后把相差这个子空间中向量的两个向量看成同一个类。

设 $M$ 是左 $R$-模，$N$ 是 $M$ 的子模。
对 $x,y\in M$，定义
\[
x\sim y \quad\Longleftrightarrow\quad x-y\in N.
\]
元素 $x$ 的等价类记为
\[
x+N=\{x+z:z\in N\},
\]
所有等价类组成的集合记为 $M/N$。

在 $M/N$ 上定义加法和 $R$ 的作用：
\[
(x+N)+(y+N)=(x+y)+N,
\qquad
r(x+N)=rx+N.
\]
同样需要检查良定义。
若 $x-x'\in N$，$y-y'\in N$，则
\[
(x+y)-(x'+y')=(x-x')+(y-y')\in N,
\]
所以加法良定义。
若 $x-x'\in N$，则
\[
rx-rx'=r(x-x')\in N,
\]
这里用到的是 $N$ 对 $R$ 的作用封闭。
因此标量作用也良定义。

\begin{definition}[商模]\label{def:quotient-module}
设 $M$ 是左 $R$-模，$N$ 是 $M$ 的子模。
按上述等价关系、加法和 $R$ 的作用得到的左 $R$-模称为 $M$ 对 $N$ 的\textbf{商模}，记为 $M/N$。
\end{definition}

\paragraph*{例：线性空间的商空间}\label{ex:quotient-vector-space}
取 $M=\mathbb R^2$，把它看成通常的实向量空间，也就是 $\mathbb R$-模。
令
\[
N=\langle(1,0)\rangle=\{(t,0):t\in\mathbb R\}.
\]
这是 $\mathbb R^2$ 的一个子模。
在商模 $\mathbb R^2/N$ 中，等价类
\[
(x,y)+N
\]
包含所有
\[
(x+t,y),\qquad t\in\mathbb R.
\]
也就是说，沿着横轴方向的差别都被看成同一个类；真正留下来的信息只有第二个坐标。

定义
\[
\rho:\mathbb R^2/N\to\mathbb R,\qquad \rho((x,y)+N)=y.
\]
这个映射是良定义的。
若
\[
(x,y)+N=(x',y')+N,
\]
则 $(x,y)-(x',y')=(x-x',y-y')\in N$。
而 $N$ 中元素的第二个坐标都是 $0$，所以 $y-y'=0$，即 $y=y'$。
因此 $\rho$ 不依赖代表元的选择。

加法和标量作用也都按第二个坐标计算：
\[
\rho(((x,y)+N)+((x',y')+N))=y+y',
\qquad
\rho(\lambda((x,y)+N))=\lambda y.
\]
所以
\[
\mathbb R^2/N\cong\mathbb R.
\]
这个例子说明，商模不是把子模简单删掉，而是把沿着子模方向的差别视为同一个类。

商模不需要额外引入"正规子模"这个概念。
原因是模的加法群本来就是阿贝尔群，而阿贝尔群中的子群自动正规；所以用差属于 $N$ 来定义等价关系时，不会遇到群论里非正规子群带来的问题。

当环 $R$ 看成自身上的左 $R$-模时，子模就是对加法封闭并且对左乘封闭的子集，也就是左理想。
如果 $R$ 是交换环，或者我们同时要求左右乘法封闭，就回到前面讲的理想。
这解释了为什么商环和商模的构造如此相似：理想本来就可以看成环作用在自身上时的特殊子模。

若 $I$ 是 $R$ 的理想，那么 $R/I$ 不只是商环，也自然是一个左 $R$-模。
作用定义为
\[
r(a+I)=ra+I.
\]
它的良定义与商模的良定义完全一样。
这个视角会在后面理解同态定理时继续出现：同态的核告诉我们哪些元素要被压成零，而商结构记录了压缩之后剩下的对象。

\subsection*{课堂问题参考解答}

\begin{exercise}\label{ex:ans-quotient-ring-two-sided}
\hypertarget{ans:quotient-ring-two-sided}{}
\textbf{商环的加法良定义只用到 $I$ 是加法子群，乘法良定义为什么必须要求 $I$ 是双边理想？}

加法良定义只需要处理
\[
(a+b)-(a'+b')=(a-a')+(b-b').
\]
如果 $a-a'\in I$，$b-b'\in I$，那么右边仍在 $I$ 中；这只用到 $I$ 是加法子群。

乘法良定义要处理
\[
ab-a'b'=(a-a')b+a'(b-b').
\]
其中 $(a-a')b$ 要求 $I$ 对右乘吸收，$a'(b-b')$ 要求 $I$ 对左乘吸收。
所以在非交换环中必须要求 $I$ 是双边理想。
如果 $R$ 是交换环，左右吸收合并成同一件事，条件看起来就没有那么明显，但本质上仍然是在保证乘法不依赖代表元。

\hyperlink{q:quotient-ring-two-sided}{返回问题}
\end{exercise}

\section{环同态第一基本定理}\label{sec:first-isomorphism-theorem-rings}

上一节的自然映射
\[
\pi:R\to R/I,\qquad \pi(a)=a+I
\]
有一个很清楚的特点：它的核正好是 $I$。
也就是说，商环 $R/I$ 是把 $I$ 中的元素全部压成零以后得到的对象。
环同态第一基本定理把这个现象反过来说：任意环同态都可以理解为先商掉它的核，再进入它的像。

\begin{theorem}[环同态第一基本定理]\label{thm:first-isomorphism-theorem-rings}
设 $\varphi:R\to S$ 是环同态。
则 $\ker\varphi$ 是 $R$ 的双边理想，并且
\[
R/\ker\varphi\cong \operatorname{im}\varphi.
\]
若记 $I=\ker\varphi$，这个同构由
\[
\psi:R/I\to\operatorname{im}\varphi,\qquad
\psi(a+I)=\varphi(a)
\]
给出。
\end{theorem}

这条定理可以先用一张交换图理解：
\[
\begin{tikzpicture}[node distance=2.8cm, baseline=(R.base)]
  \node (R) {$R$};
  \node[right=of R] (Im) {$\operatorname{im}\varphi$};
  \node[below=of R] (Q) {$R/\ker\varphi$};
  \draw[->] (R) -- node[above] {$\varphi$} (Im);
  \draw[->] (R) -- node[left] {$\pi$} (Q);
  \draw[->] (Q) -- node[below right] {$\psi$} (Im);
\end{tikzpicture}
\]
其中 $\pi(a)=a+\ker\varphi$ 是自然投影，$\psi(a+\ker\varphi)=\varphi(a)$ 是定理构造出的同构。
图的意思是：从 $R$ 直接沿上方的 $\varphi$ 走到 $\operatorname{im}\varphi$，和先沿下方的 $\pi$ 走到商环，再沿 $\psi$ 走到像，结果相同。
用公式写就是
\[
\varphi=\psi\circ\pi.
\]
这句话是第一基本定理的核心读法：同态 $\varphi$ 的"混淆"全部由核负责；商掉核以后，剩下的部分与像完全同构。

\paragraph*{证明}
先证明 $\ker\varphi$ 是 $R$ 的双边理想。
因为环同态保持零元，所以
\[
\varphi(0_R)=0_S,
\]
从而 $0_R\in\ker\varphi$，核非空。
若 $a,b\in\ker\varphi$，则
\[
\varphi(a-b)=\varphi(a)-\varphi(b)=0_S-0_S=0_S,
\]
所以 $a-b\in\ker\varphi$。
这说明 $\ker\varphi$ 对差封闭。

还要检查它能吸收 $R$ 中任意元素的左右乘法。
任取 $r\in R$，$a\in\ker\varphi$。
由于 $\varphi(a)=0_S$，有
\[
\varphi(ra)=\varphi(r)\varphi(a)=\varphi(r)0_S=0_S,
\]
所以 $ra\in\ker\varphi$。
同理
\[
\varphi(ar)=\varphi(a)\varphi(r)=0_S\varphi(r)=0_S,
\]
所以 $ar\in\ker\varphi$。
左右两边都要验证，是因为这里的环不一定交换。
因此 $\ker\varphi$ 是双边理想。

下面构造从商环到像的映射。
记
\[
I=\ker\varphi.
\]
商环中的元素具有形式 $a+I$。
定义
\[
\psi:R/I\to\operatorname{im}\varphi,\qquad \psi(a+I)=\varphi(a).
\]
这里要注意，$\psi(a+I)$ 映成的是 $\varphi(a)$，不是 $a$。
映射 $\varphi$ 的输入是 $R$ 中的元素，而 $\psi$ 的输入是商环中的等价类。
定义式右边的 $\varphi(a)$ 按定义属于 $\operatorname{im}\varphi$，所以值域没有跑出去。

先检查良定义。
若 $a+I=b+I$，则 $a-b\in I=\ker\varphi$。
因此
\[
0_S=\varphi(a-b)=\varphi(a)-\varphi(b),
\]
从而 $\varphi(a)=\varphi(b)$。
所以 $\psi(a+I)$ 不依赖代表元的选择。

再检查 $\psi$ 保持环运算。
对加法，有
\[
\psi((a+I)+(b+I))
=\psi((a+b)+I)
=\varphi(a+b)
=\varphi(a)+\varphi(b)
=\psi(a+I)+\psi(b+I).
\]
对乘法，有
\[
\psi((a+I)(b+I))
=\psi(ab+I)
=\varphi(ab)
=\varphi(a)\varphi(b)
=\psi(a+I)\psi(b+I).
\]
按照本讲义采用的环同态约定，这里不额外验证保持单位元；需要保持单位元时，只需在同态定义中额外加入相应条件。

接着证明 $\psi$ 是单射。
若
\[
\psi(a+I)=\psi(b+I),
\]
则
\[
\varphi(a)=\varphi(b).
\]
移项得到
\[
\varphi(a-b)=0_S,
\]
所以 $a-b\in\ker\varphi=I$。
这说明 $a$ 和 $b$ 在商环中属于同一个等价类，即
\[
a+I=b+I.
\]
注意，这里并不是推出 $a=b$。
因为 $\varphi$ 不一定是单射，我们只能推出 $a-b$ 落在核里；但这已经足够说明两个等价类相等。

\hypertarget{q:first-isomorphism-not-a-equals-b}{}
\medskip
\noindent\textbf{思考：为什么从 $\psi(a+I)=\psi(b+I)$ 只能推出 $a+I=b+I$，而不能推出 $a=b$？}（\hyperlink{ans:first-isomorphism-not-a-equals-b}{跳转到解答}。）
\medskip

最后证明 $\psi$ 是满射。
任取 $y\in\operatorname{im}\varphi$。
按像的定义，存在 $a\in R$，使得
\[
y=\varphi(a).
\]
于是
\[
\psi(a+I)=\varphi(a)=y.
\]
这就为任意像中的元素找到了商环中的原像。
因此 $\psi$ 是双射环同态，也就是环同构。
定理得证。 $\square$

\subsection*{课堂问题参考解答}

\begin{exercise}\label{ex:ans-first-isomorphism-not-a-equals-b}
\hypertarget{ans:first-isomorphism-not-a-equals-b}{}
\textbf{为什么从 $\psi(a+I)=\psi(b+I)$ 只能推出 $a+I=b+I$，而不能推出 $a=b$？}

等式 $\psi(a+I)=\psi(b+I)$ 的意思是
\[
\varphi(a)=\varphi(b).
\]
这只说明
\[
\varphi(a-b)=0,
\]
也就是 $a-b\in\ker\varphi=I$。
商环 $R/I$ 正是把相差一个 $I$ 中元素的两个元素看成同一个类，所以我们能推出的是
\[
a+I=b+I.
\]
如果核不为零，原环中确实可能有 $a\neq b$ 但 $\varphi(a)=\varphi(b)$。
第一基本定理的要点正是：把核商掉以后，这种混淆就消失了。

\hyperlink{q:first-isomorphism-not-a-equals-b}{返回问题}
\end{exercise}

\section{后续同构定理}\label{sec:further-isomorphism-theorems-rings}

在第一基本定理中，我们研究的是一个同态的核，以及把核商掉以后得到的结构。
但在一个环 $R$ 中，常常会同时出现多个理想。
多个理想之间一旦发生关系，就会衍生出更复杂的商结构。

最简单的情况是两个理想有包含关系。
设 $I,J$ 都是 $R$ 的理想，并且
\[
I\subseteq J.
\]
这时有两个自然问题：
一方面，$J/I$ 能不能看成商环 $R/I$ 中的理想；
另一方面，$R/I$、$R/J$ 与 $J/I$ 之间究竟有什么关系。
这条线索导向第三基本定理。

另一种情况是 $I$ 和 $J$ 没有包含关系。
这时不能直接写出 $J/I$ 或 $I/J$，于是自然要看两个由 $I,J$ 共同生成的对象：
\[
I+J,\qquad I\cap J.
\]
前者记录把两个理想合在一起以后得到的部分，后者记录二者的重叠部分。
这条线索导向第二基本定理。

\subsection{满同态下的理想对应}\label{sec:ideal-correspondence-quotient}

在正式讨论第三、第二基本定理之前，先看一个常用的中间结论。
它说明满同态会把包含核的理想，与目标环中的理想联系起来。

\begin{theorem}[满同态下的理想对应]\label{thm:ideal-correspondence-quotient}
设 $\varphi:R\to S$ 是满环同态，$J=\ker\varphi$。

若 $R'$ 是 $R$ 的双边理想且 $J\subseteq R'$，则 $S'=\varphi(R')$ 是 $S$ 的双边理想。
若 $S'$ 是 $S$ 的双边理想，则 $\varphi^{-1}(S')$ 是 $R$ 的双边理想，并且包含 $J$。

进一步，若 $R'$ 与 $S'=\varphi(R')$ 对应，则有自然同构
\[
R/R'\cong S/S'.
\]
\end{theorem}

\paragraph*{证明}
先证明像的方向。
设 $R'$ 是 $R$ 的双边理想，且 $J\subseteq R'$。
令
\[
S'=\varphi(R')=\{\varphi(r'):r'\in R'\}.
\]
因为 $0_R\in R'$，所以
\[
0_S=\varphi(0_R)\in S',
\]
从而 $S'$ 非空。
若 $s_1,s_2\in S'$，则存在 $r_1,r_2\in R'$，使得
\[
s_1=\varphi(r_1),\qquad s_2=\varphi(r_2).
\]
于是
\[
s_1-s_2=\varphi(r_1)-\varphi(r_2)=\varphi(r_1-r_2).
\]
由于 $R'$ 对差封闭，$r_1-r_2\in R'$，所以 $s_1-s_2\in S'$。
这说明 $S'$ 对差封闭。

接着检查吸收性。
任取 $s_1\in S'$，任取 $s\in S$。
写 $s_1=\varphi(r_1)$，其中 $r_1\in R'$。
因为 $\varphi$ 是满射，存在 $r\in R$，使得
\[
s=\varphi(r).
\]
于是
\[
ss_1=\varphi(r)\varphi(r_1)=\varphi(rr_1),
\qquad
s_1s=\varphi(r_1)\varphi(r)=\varphi(r_1r).
\]
由于 $R'$ 是双边理想，$rr_1,r_1r\in R'$。
因此 $ss_1,s_1s\in S'$。
所以 $S'$ 是 $S$ 的双边理想。
注意，满射正是在这里使用的：为了处理任意 $S$ 中元素 $s$，必须把它写成某个 $R$ 中元素的像。

再证明原像的方向。
设 $S'$ 是 $S$ 的双边理想。
记
\[
T=\varphi^{-1}(S')=\{r\in R:\varphi(r)\in S'\}.
\]
因为 $0_S\in S'$ 且 $\varphi(0_R)=0_S$，所以 $0_R\in T$。
若 $x,y\in T$，则 $\varphi(x),\varphi(y)\in S'$，从而
\[
\varphi(x-y)=\varphi(x)-\varphi(y)\in S',
\]
所以 $x-y\in T$。
这说明 $T$ 对差封闭。

任取 $x\in T$，$r\in R$。
由 $x\in T$ 可知 $\varphi(x)\in S'$。
于是
\[
\varphi(rx)=\varphi(r)\varphi(x)\in S',
\qquad
\varphi(xr)=\varphi(x)\varphi(r)\in S',
\]
因为 $S'$ 是 $S$ 的双边理想。
所以 $rx,xr\in T$。
因此 $T$ 是 $R$ 的双边理想。
这一方向没有用到 $\varphi$ 的满射性。
此外，若 $a\in J=\ker\varphi$，则
\[
\varphi(a)=0_S\in S',
\]
所以 $a\in T$。
这说明 $J\subseteq \varphi^{-1}(S')$。

最后证明商环之间的同构。
设 $R'$ 是包含 $J$ 的双边理想，并令 $S'=\varphi(R')$。
定义
\[
\Psi:R/R'\to S/S',
\qquad
\Psi(a+R')=\varphi(a)+S'.
\]
这个定义的值确实落在 $S/S'$ 中。
下面先检查良定义。
若
\[
a+R'=b+R',
\]
则 $a-b\in R'$。
因此
\[
\varphi(a)-\varphi(b)=\varphi(a-b)\in \varphi(R')=S',
\]
所以
\[
\varphi(a)+S'=\varphi(b)+S'.
\]
这说明 $\Psi$ 不依赖代表元的选择。

再检查 $\Psi$ 保持环运算。
对加法，有
\[
\Psi((a+R')+(b+R'))
=\Psi((a+b)+R')
=\varphi(a+b)+S'
\]
\[
=(\varphi(a)+\varphi(b))+S'
=(\varphi(a)+S')+(\varphi(b)+S')
=\Psi(a+R')+\Psi(b+R').
\]
对乘法，有
\[
\Psi((a+R')(b+R'))
=\Psi(ab+R')
=\varphi(ab)+S'
\]
\[
=\varphi(a)\varphi(b)+S'
=(\varphi(a)+S')(\varphi(b)+S')
=\Psi(a+R')\Psi(b+R').
\]
所以 $\Psi$ 是环同态。

证明 $\Psi$ 是满射。
任取 $s+S'\in S/S'$。
因为 $\varphi$ 是满射，存在 $a\in R$，使得 $s=\varphi(a)$。
于是
\[
\Psi(a+R')=\varphi(a)+S'=s+S'.
\]
所以 $\Psi$ 满射。

证明 $\Psi$ 是单射。
只需说明 $\ker\Psi$ 是 $R/R'$ 的零元。
若
\[
\Psi(a+R')=S',
\]
则
\[
\varphi(a)+S'=S',
\]
也就是 $\varphi(a)\in S'=\varphi(R')$。
因此存在 $r'\in R'$，使得
\[
\varphi(a)=\varphi(r').
\]
于是
\[
\varphi(a-r')=0_S,
\]
所以 $a-r'\in\ker\varphi=J$。
又因为 $J\subseteq R'$，且 $r'\in R'$，可得
\[
a=(a-r')+r'\in R'.
\]
这说明 $a+R'=R'$。
因此 $\ker\Psi$ 只有零元，$\Psi$ 是单射。
综上，
\[
R/R'\cong S/S'.
\]
定理证毕。 $\square$

这里也可以看出满射条件的重要性。
例如从 $\mathbb Z$ 到 $\mathbb Q$ 的包含同态不是满同态，$\mathbb Z$ 中的理想在 $\mathbb Q$ 中的像不一定是 $\mathbb Q$ 的理想。
比如 $2\mathbb Z$ 在 $\mathbb Q$ 中并不能吸收任意有理数乘法，所以不是 $\mathbb Q$ 的理想。

\hypertarget{q:ideal-correspondence-surjective}{}
\medskip
\noindent\textbf{思考：为什么证明 $\varphi(R')$ 是 $S$ 的理想时必须用到 $\varphi$ 满射？}（\hyperlink{ans:ideal-correspondence-surjective}{跳转到解答}。）
\medskip

\subsection{第三基本定理：从分解图看商中再商}\label{sec:third-isomorphism-theorem-diagram}

现在回到包含关系的情况。
设 $\varphi:R\to S$ 是满环同态，$J=\ker\varphi$。
如果 $I$ 是 $R$ 的双边理想，并且
\[
I\subseteq J,
\]
那么 $\varphi$ 会经过商环 $R/I$ 下降为一个同态。
这件事可以用下面的图来组织：
\[
\begin{tikzpicture}[node distance=2.7cm and 3.4cm, baseline=(R.base)]
  \node (R) {$R$};
  \node[right=of R] (RJ) {$R/J$};
  \node[below=of R] (RI) {$R/I$};
  \node[right=of RI] (S) {$S$};
  \node[below=of RI] (Q) {$(R/I)/(J/I)$};

  \draw[->] (R) -- node[above] {$\pi_J$} (RJ);
  \draw[->] (R) -- node[left] {$\pi_I$} (RI);
  \draw[->] (R) -- node[above right] {$\varphi$} (S);
  \draw[->] (RJ) -- node[right] {$\varphi_J$} (S);
  \draw[->] (RI) -- node[above] {$\bar\varphi$} (S);
  \draw[->] (RI) -- node[left] {$\pi_{J/I}$} (Q);
  \draw[->] (Q) -- node[below right] {$\cong$} (S);
\end{tikzpicture}
\]
这里 $\pi_I$ 与 $\pi_J$ 都是自然映射。
由第一基本定理，$\varphi_J:R/J\to S$ 是同构。
图中真正需要新构造的是
\[
\bar\varphi:R/I\to S,\qquad \bar\varphi(a+I)=\varphi(a).
\]
如果这个映射构造成功，并且它的核是 $J/I$，那么再对 $\bar\varphi$ 使用第一基本定理，就会得到
\[
(R/I)/(J/I)\cong S.
\]
又因为 $R/J\cong S$，于是得到
\[
(R/I)/(J/I)\cong R/J.
\]

先说明 $J/I$ 可以看成 $R/I$ 的双边理想。
它的元素形式是
\[
j+I,\qquad j\in J.
\]
因为 $J$ 对差封闭，若 $j_1,j_2\in J$，则
\[
(j_1+I)-(j_2+I)=(j_1-j_2)+I\in J/I.
\]
又任取 $r\in R$，$j\in J$，在 $R/I$ 中有
\[
(r+I)(j+I)=rj+I,\qquad (j+I)(r+I)=jr+I.
\]
由于 $J$ 是 $R$ 的双边理想，$rj,jr\in J$，所以这两个乘积仍在 $J/I$ 中。
因此 $J/I$ 是 $R/I$ 的双边理想。

下面构造 $\bar\varphi$。
定义
\[
\bar\varphi:R/I\to S,\qquad \bar\varphi(a+I)=\varphi(a).
\]
先检查良定义。
若 $a+I=b+I$，则 $a-b\in I$。
因为 $I\subseteq J=\ker\varphi$，所以 $a-b\in\ker\varphi$，从而
\[
0_S=\varphi(a-b)=\varphi(a)-\varphi(b).
\]
因此 $\varphi(a)=\varphi(b)$。
这说明 $\bar\varphi$ 不依赖代表元选择。

再检查它是环同态。
对加法，
\[
\bar\varphi((a+I)+(b+I))
=\bar\varphi((a+b)+I)
=\varphi(a+b)
=\varphi(a)+\varphi(b)
=\bar\varphi(a+I)+\bar\varphi(b+I).
\]
对乘法，
\[
\bar\varphi((a+I)(b+I))
=\bar\varphi(ab+I)
=\varphi(ab)
=\varphi(a)\varphi(b)
=\bar\varphi(a+I)\bar\varphi(b+I).
\]
所以 $\bar\varphi$ 是环同态。

它还是满射。
任取 $s\in S$，由于 $\varphi$ 是满射，存在 $a\in R$，使得 $s=\varphi(a)$。
于是
\[
\bar\varphi(a+I)=\varphi(a)=s.
\]

最后计算 $\bar\varphi$ 的核。
若 $a+I\in\ker\bar\varphi$，则
\[
\bar\varphi(a+I)=0_S,
\]
即 $\varphi(a)=0_S$。
所以 $a\in\ker\varphi=J$，从而
\[
a+I\in J/I.
\]
反过来，若 $a+I\in J/I$，则 $a\in J=\ker\varphi$，于是
\[
\bar\varphi(a+I)=\varphi(a)=0_S.
\]
所以 $a+I\in\ker\bar\varphi$。
因此
\[
\ker\bar\varphi=J/I.
\]

\hypertarget{q:third-iso-kernel}{}
\medskip
\noindent\textbf{思考：为什么 $\ker\bar\varphi$ 是 $J/I$，而不是 $J$？}（\hyperlink{ans:third-iso-kernel}{跳转到解答}。）
\medskip

对 $\bar\varphi$ 使用第一基本定理，得到
\[
(R/I)/(J/I)\cong S.
\]
另一方面，对原来的满同态 $\varphi:R\to S$ 使用第一基本定理，得到
\[
R/J\cong S.
\]
于是
\[
(R/I)/(J/I)\cong R/J.
\]
这说明：如果 $I\subseteq J$，先商掉 $I$，再在商环 $R/I$ 中商掉 $J/I$，和一开始直接商掉 $J$ 得到的是同构的结果。

\subsection{第三基本定理的常规形式}\label{sec:third-isomorphism-theorem-rings}

把上面的讨论抽离出来，就得到第三基本定理的常规表述。

\begin{theorem}[环同态第三基本定理]\label{thm:third-isomorphism-theorem-rings}
设 $R$ 是环，$I,J$ 是 $R$ 的双边理想，且 $I\subseteq J$。
则 $J/I$ 是 $R/I$ 的双边理想，并且
\[
(R/I)/(J/I)\cong R/J.
\]
\end{theorem}

\paragraph*{证明}
前面已经说明过，$J/I$ 是 $R/I$ 的双边理想。
现在直接构造一个从 $R/I$ 到 $R/J$ 的满同态：
\[
\Phi:R/I\to R/J,\qquad \Phi(r+I)=r+J.
\]
先检查良定义。
若 $r_1+I=r_2+I$，则 $r_1-r_2\in I$。
由于 $I\subseteq J$，有 $r_1-r_2\in J$，所以
\[
r_1+J=r_2+J.
\]
因此 $\Phi$ 良定义。

它是满射，因为任取 $r+J\in R/J$，都有
\[
\Phi(r+I)=r+J.
\]
它保持加法：
\[
\Phi((r_1+I)+(r_2+I))
=\Phi((r_1+r_2)+I)
=(r_1+r_2)+J
\]
\[
=(r_1+J)+(r_2+J)
=\Phi(r_1+I)+\Phi(r_2+I).
\]
它保持乘法：
\[
\Phi((r_1+I)(r_2+I))
=\Phi(r_1r_2+I)
=r_1r_2+J
\]
\[
=(r_1+J)(r_2+J)
=\Phi(r_1+I)\Phi(r_2+I).
\]
所以 $\Phi$ 是满环同态。

再计算核。
一个元素 $r+I$ 落在 $\ker\Phi$ 中，当且仅当
\[
\Phi(r+I)=J,
\]
也就是
\[
r+J=J.
\]
这等价于 $r\in J$。
因此
\[
\ker\Phi=\{r+I:r\in J\}=J/I.
\]
由第一基本定理，
\[
(R/I)/(J/I)\cong R/J.
\]
定理证毕。 $\square$

\subsection{第二基本定理：相加与相交}\label{sec:second-isomorphism-theorem-rings}

第三基本定理处理的是 $I\subseteq J$ 的情形。
如果两个理想没有包含关系，就不能直接写出 $J/I$ 或 $I/J$。
这时自然出现两个新对象：一个是把两边都收进来的 $I+J$，另一个是两边共有的 $I\cap J$。
第二基本定理正是在比较这两个对象给出的商结构。

本校教材中常用的是稍微更一般的版本：一个子环加一个双边理想。

\begin{theorem}[环同态第二基本定理]\label{thm:second-isomorphism-theorem-rings}
设 $R$ 是环，$A$ 是 $R$ 的子环，$I$ 是 $R$ 的双边理想。
则 $A+I$ 是 $R$ 的子环，$A\cap I$ 是 $A$ 的双边理想，并且
\[
A/(A\cap I)\cong (A+I)/I.
\]
\end{theorem}

\paragraph*{证明}
先证明 $A+I$ 是 $R$ 的子环。
这里
\[
A+I=\{a+x:a\in A,\ x\in I\}.
\]
因为 $0_R=0_R+0_R$，且 $0_R\in A$、$0_R\in I$，所以 $A+I$ 非空。

任取 $a_1+x_1,a_2+x_2\in A+I$，其中 $a_1,a_2\in A$，$x_1,x_2\in I$。
则
\[
(a_1+x_1)-(a_2+x_2)=(a_1-a_2)+(x_1-x_2).
\]
由于 $A$ 是子环，$a_1-a_2\in A$；
由于 $I$ 是理想，$x_1-x_2\in I$。
所以差仍属于 $A+I$。

再检查乘法封闭。
有
\[
(a_1+x_1)(a_2+x_2)
=a_1a_2+a_1x_2+x_1a_2+x_1x_2.
\]
其中 $a_1a_2\in A$。
后三项都属于 $I$：$a_1x_2\in I$、$x_1a_2\in I$、$x_1x_2\in I$。
因此
\[
a_1x_2+x_1a_2+x_1x_2\in I.
\]
整个乘积仍然可以写成一个 $A$ 中元素加一个 $I$ 中元素，所以属于 $A+I$。
因此 $A+I$ 是 $R$ 的子环。

接着证明 $A\cap I$ 是 $A$ 的双边理想。
它非空，因为 $0_R\in A\cap I$。
若 $x,y\in A\cap I$，则 $x-y\in A$ 且 $x-y\in I$，所以 $x-y\in A\cap I$。
任取 $a\in A$，$x\in A\cap I$。
因为 $A$ 是子环，$ax,xa\in A$；
因为 $I$ 是 $R$ 的双边理想，$ax,xa\in I$。
所以
\[
ax,xa\in A\cap I.
\]
这说明 $A\cap I$ 是 $A$ 的双边理想。

现在构造同态
\[
\theta:A\to (A+I)/I,\qquad \theta(a)=a+I.
\]
它保持加法：
\[
\theta(a+b)=(a+b)+I=(a+I)+(b+I)=\theta(a)+\theta(b).
\]
它保持乘法：
\[
\theta(ab)=ab+I=(a+I)(b+I)=\theta(a)\theta(b).
\]

证明 $\theta$ 是满射。
任取 $(A+I)/I$ 中的元素 $z+I$，其中 $z\in A+I$。
按 $A+I$ 的定义，存在 $a\in A$、$i\in I$，使得
\[
z=a+i.
\]
于是
\[
z+I=(a+i)+I=a+I=\theta(a).
\]
所以 $\theta$ 满射。

\hypertarget{q:second-iso-surjective}{}
\medskip
\noindent\textbf{思考：为什么 $a+i+I=a+I$ 说明 $\theta:A\to(A+I)/I$ 是满射？}（\hyperlink{ans:second-iso-surjective}{跳转到解答}。）
\medskip

最后计算核。
若 $a\in\ker\theta$，则
\[
\theta(a)=I,
\]
也就是
\[
a+I=I.
\]
这等价于 $a\in I$。
而 $\theta$ 的定义域是 $A$，所以 $a\in A$。
因此 $a\in A\cap I$。
反过来，若 $a\in A\cap I$，则 $a\in I$，所以
\[
\theta(a)=a+I=I,
\]
即 $a\in\ker\theta$。
因此
\[
\ker\theta=A\cap I.
\]
由第一基本定理可得
\[
A/(A\cap I)\cong (A+I)/I.
\]
定理证毕。 $\square$

如果把 $A$ 换成另一个双边理想 $J$，就得到常见的理想版本：
\[
J/(J\cap I)\cong (J+I)/I.
\]
这正对应前面说的第二种情况：当两个理想没有包含关系时，我们通过相加 $J+I$ 和相交 $J\cap I$ 来描述它们之间的商结构。

\subsection*{课堂问题参考解答}

\begin{exercise}\label{ex:ans-ideal-correspondence-surjective}
\hypertarget{ans:ideal-correspondence-surjective}{}
\textbf{为什么证明 $\varphi(R')$ 是 $S$ 的理想时必须用到 $\varphi$ 满射？}

要证明 $\varphi(R')$ 是 $S$ 的理想，吸收性要求对任意 $s\in S$ 和任意 $s_1\in\varphi(R')$，都有
\[
ss_1,\ s_1s\in\varphi(R').
\]
若 $s_1=\varphi(r_1)$，其中 $r_1\in R'$，那么为了把 $ss_1$ 写成某个 $R'$ 中元素的像，必须先把任意的 $s\in S$ 写成 $\varphi(r)$。
这正是满射条件提供的。
如果 $\varphi$ 不是满射，目标环中可能有一些元素根本不是任何 $R$ 中元素的像，吸收性就无法从 $R'$ 的吸收性推出。

\hyperlink{q:ideal-correspondence-surjective}{返回问题}
\end{exercise}

\begin{exercise}\label{ex:ans-third-iso-kernel}
\hypertarget{ans:third-iso-kernel}{}
\textbf{为什么 $\ker\bar\varphi$ 是 $J/I$，而不是 $J$？}

核必须是定义域中的子集。
这里 $\bar\varphi$ 的定义域是 $R/I$，所以核中的元素不是 $R$ 中的元素 $a$，而是商环中的等价类 $a+I$。

由定义，
\[
\bar\varphi(a+I)=0_S
\]
等价于
\[
\varphi(a)=0_S,
\]
也就是 $a\in J$。
因此核中的元素正是所有形如 $a+I$ 且 $a\in J$ 的等价类：
\[
\ker\bar\varphi=\{a+I:a\in J\}=J/I.
\]
所以这里不能写成 $J$；$J$ 是 $R$ 中的理想，而 $J/I$ 才是 $R/I$ 中的理想。

\hyperlink{q:third-iso-kernel}{返回问题}
\end{exercise}

\begin{exercise}\label{ex:ans-second-iso-surjective}
\hypertarget{ans:second-iso-surjective}{}
\textbf{为什么 $a+i+I=a+I$ 说明 $\theta:A\to(A+I)/I$ 是满射？}

满射要证明的是：$(A+I)/I$ 中每个等价类都能被 $\theta$ 命中。
任取 $z+I\in (A+I)/I$。
因为 $z\in A+I$，所以可以写成
\[
z=a+i,\qquad a\in A,\ i\in I.
\]
在模掉 $I$ 的商环里，$i$ 被看成零，因此
\[
z+I=(a+i)+I=a+I.
\]
而 $a\in A$，所以
\[
a+I=\theta(a).
\]
这就为任意 $z+I$ 找到了来自 $A$ 的原像 $a$，因此 $\theta$ 是满射。

\hyperlink{q:second-iso-surjective}{返回问题}
\end{exercise}
